\documentclass{article}

% Language setting
% Replace `english' with e.g. `spanish' to change the document language
\usepackage[english]{babel}


% Set page size and margins
% Replace `letterpaper' with `a4paper' for UK/EU standard size
\usepackage[letterpaper,top=2cm,bottom=2cm,left=3cm,right=3cm,marginparwidth=1.75cm]{geometry}

% Useful packages
\usepackage{amsmath}
\usepackage{graphicx}
\usepackage[colorlinks=true, allcolors=blue]{hyperref}
\usepackage{graphicx}
\usepackage{subcaption}
\usepackage{amssymb}
\usepackage{xcolor}


% Simple read marker used by the status boxes
\newcommand{\ReadMark}[1]{%
  \ifnum#1=1
    \textcolor{green!50!black}{\textbf{[read]}}
  \else
    \textcolor{red!70!black}{\textbf{[not read]}}
  \fi
}

% Summary status for 2 authors
\newcommand{\AllReadTwo}[2]{%
  \ifnum\numexpr#1+#2\relax=2
    \textcolor{green!50!black}{\textbf{Both authors have read this part.}}
  \else
    \textcolor{orange!80!black}{\textbf{Waiting for confirmation from both authors.}}
  \fi
}

% #1 = section/part name
% #2 = author 1 name, #3 = 0/1
% #4 = author 2 name, #5 = 0/1
\newcommand{\ReadStatusTwo}[5]{%
  \par\smallskip
  \noindent
  \fcolorbox{black!20}{gray!8}{%
    \parbox{0.97\linewidth}{%
      \textbf{Reading status: #1}\\[3pt]
      #2~\ReadMark{#3}\qquad
      #4~\ReadMark{#5}\\[4pt]
      \AllReadTwo{#3}{#5}
    }%
  }%
  \par\smallskip
}

% Default subsection status:
% change the two 0 values below to 1 when Yulian or Michael confirms reading
\newcommand{\SubsectionStatus}[1]{%
  \ReadStatusTwo{#1}{Yulian}{0}{Michael}{0}
}

\title{Two-Scale Gillespie Simulation on Community-Structured Networks:\\
Complexity Reduction, Analytical Validation, and Accuracy Trade-offs}
\author{Yulian \and Michael}

\title{Notes on Macro layer of Two-Scale Stochastic Simulator}

\begin{document}
\maketitle

% \begin{abstract}
% Your abstract.
% \end{abstract}

\tableofcontents

\begin{abstract}
    We study a two-scale continuous-time stochastic simulation framework for spreading processes on community-structured networks. The method preserves exact event-driven CTMC dynamics within communities while replacing explicit cross-community edges by macro-level stochastic hazards driven by aggregate community states. We formalize the full-network Gillespie CTMC as a reference process and decompose the approximation error of the two-scale simulator into rate abstraction, import placement, and numerical synchronization components. On analytically tractable clique-bridge-leaf motifs, we show that aggregate mean-field hazards are conditional means of the exact bridge hazards under exchangeability assumptions, but may nevertheless create premature export timing through Jensen-type survival effects. We extend the analysis to chains of cliques, where local timing shifts accumulate across successive community activations. Numerical experiments compare arrival-time distributions and runtime against full-network simulation, identifying regimes in which the two-scale approximation is accurate and regimes in which narrow boundary bottlenecks induce systematic bias.
\end{abstract}

\section{Introduction}
\ReadStatusTwo{}{Yulian}{0}{Michael}{0}
The modeling of epidemic spreading on complex networks has become a cornerstone of modern computational epidemiology, providing critical insights for public health decision-making. However, the simulation of contagion processes on a national or global scale (involving $N > 10^7$ agents) presents a fundamental challenge: balancing the computational efficiency required for massive scenario analysis with the granular accuracy needed to capture stochastic effects, particularly during the early stages of an outbreak.


\subsection{Motivation: The Data Asymmetry Problem}
\ReadStatusTwo{}{Yulian}{1}{Michael}{0}
A primary bottleneck in constructing realistic Agent-Based Models (ABMs) is the limited availability of reliable data on the exact topology of microscopic contact networks of a large size. In a densely populated city, it is practically impossible to observe and validate every individual interaction. As a result, models that depend on a detailed synthetic reconstruction of person-to-person contacts may inherit substantial uncertainty from assumptions that are difficult to verify empirically.

In contrast, information at a coarser spatial scale is often much more accessible. Aggregated mobility data, such as commuter flows between cities, air travel volumes, and inter-regional movement patterns, can frequently be estimated from census records, transportation data, and related mobility sources. This type of information already underlies metapopulation-style epidemic models in which spatially separated populations are coupled through mobility or transportation structure \cite{colizza2008,balcan2010modeling}. This creates a natural situation of \textit{data asymmetry}: flows \textit{between} communities may be known at least approximately, whereas the internal contact structure \textit{within} communities is usually only partially observed.

This discrepancy motivates a modeling strategy that is aligned with the resolution of the available data. Rather than attempting to specify an unverifiable global microscopic network, it is natural to represent inter-community transmission through observable or estimable mobility patterns while treating local interactions stochastically. A separate but related bottleneck is computational complexity: even when exact Gillespie-style simulation provides an appropriate microscopic reference, large agent-based systems remain expensive to simulate because the process is event-driven and requires repeated updates of a large number of possible interactions.

\subsection{The Two-Scale Network Abstraction}
\ReadStatusTwo{}{Yulian}{1}{Michael}{0}

To address this challenge, we propose a \textbf{Two-Scale Network} approach based on structural decomposition. The central idea is to replace an intractable global contact graph by a hierarchical community-structured system that is close in spirit to metapopulation and multiscale epidemic models \cite{colizza2008,balcan2010modeling}.
The proposed abstraction includes:

\begin{itemize}
    \item \textbf{Macro-scale (The Meta-graph):} Cities, towns, or administrative districts are abstracted into \textbf{meta-nodes}. The connections between them are defined not by individual edges, but by weighted links representing mobility flows (e.g., passenger traffic intensity). This allows inter-community transmission to be represented through a graph of communities driven by mobility structure rather than by a full set of microscopic cross-community contacts.
    
    \item \textbf{Micro-scale (The Intra-node Dynamics):} Each meta-node is treated as a local stochastic subsystem whose intra-community dynamics are modeled separately from the macro graph. In the present work, this means that each community is represented at the micro layer by a CTMC simulated in the Gillespie algorithm, while its aggregate state contributes to the macro layer through quantities relevant for inter-community transmission.
\end{itemize}

This abstraction allows us to model the system as a set of weakly coupled stochastic oscillators, where the internal state of a meta-node evolves rapidly, while the coupling (transmission between nodes) occurs as rare, discrete events driven by the mobility weights.

\subsection{Limitations of Existing Approaches}
\ReadStatusTwo{}{Yulian}{1}{Michael}{0}

Metapopulation and multiscale epidemic models are well established in the literature. Foundational work by Colizza, Vespignani, and coauthors represents spatially separated populations as coupled units connected through transportation or mobility networks \cite{colizza2008,balcan2010modeling}. More recent large-scale modeling frameworks and multiscale agent-based approaches retain this general logic while incorporating richer local dynamics and higher-resolution data \cite{kerr2021}. These approaches are valuable because they make large-scale simulation feasible and align naturally with observable movement patterns. At the same time, a common implementation strategy in coarse-grained settings is to advance the coupled system through fixed synchronization intervals.

This time-driven viewpoint is often adequate for aggregate prevalence trajectories, but it becomes less satisfactory when inter-community transmission is driven by rare stochastic arrivals. When the import of a single infectious individual may determine whether a susceptible community experiences local fade-out or a sustained outbreak, fixed time steps can blur the distinction between no arrival, one arrival, and several arrivals within the same interval. In this sense, the synchronization scheme itself may contribute to approximation error in the representation of inter-community transmission.

From the opposite direction, the exact Stochastic Simulation Algorithm (SSA) of Gillespie provides a continuous-time framework in which event times are generated directly from current propensities rather than imposed by an external time grid \cite{gillespie1977exact}. This makes Gillespie-style simulation an appropriate microscopic reference for rare-event dynamics. There is inverse approach -  to represent reactions through independent unit-rate Poisson processes with internal times given by integrated propensity functions \cite{anderson2007}. Related hazard-integral formulations further clarify how continuous-time event timing can be handled under time-varying rates and changing contact structure \cite{vestergaard2015}. However, a direct application of exact SSA to a very large community-structured population remains computationally demanding, because the simulation must repeatedly maintain and update a large collection of possible interactions or channels.

The present work is motivated by this gap. Aggregate metapopulation models may be too coarse when rare import events play an important role, whereas full microscopic SSA models may become too expensive at large scale. This motivates a two-scale formulation that preserves a continuous-time stochastic description at the micro layer while introducing an abstracted macro layer for inter-community transmission.


\subsection{Research Objective}
\ReadStatusTwo{}{Yulian}{1}{Michael}{0}

This paper introduces a novel \textbf{Two-Scale Stochastic Simulation framework} that bridges the gap between data avaliability, stochastic fidelity and computational feasibility. The basic idea is to represent communities as meta-nodes at the macro layer, while retaining a microscopic continuous-time stochastic description for intra-community dynamics at the micro layer. Inter-community transmission is then modeled through mobility-weighted stochastic hazards whose timing is handled in continuous time rather than through fixed synchronization intervals. Our specific contributions are:

\begin{enumerate}
    \item A formal definition of the \textbf{Two-Scale architecture}, where communities are modeled as meta-nodes with internal stochastic kinetics, linked by mobility-weighted hazards. The meta-nodes interact through a macro graph while preserving stochastic local dynamics within each community.
    \item The development of a synchronization mechanism with the logic of Gillespie framework, based on the \textbf{Hazard Integral} (Cumulative Hazard), which allows for the exact timing of inter-node transmission events without fixed time steps. This ensures that the model preserves the "memory" of risk accumulation, correctly simulating the probability of infection transfer even through weak mobility channels.
    \item \textbf{Analysis of the accuracy--efficiency trade-off} introduced by the two-scale abstraction. 
    We distinguish the exact part of the simulator, namely the local CTMC/Gillespie dynamics inside communities, from the approximate part, namely the macro-level representation of inter-community transmission. 
    In particular, we study how replacing boundary-node states by aggregate community fractions can create interface-level timing bias relative to a full-network Gillespie reference.
    \item \textbf{Analytical and numerical interface-error benchmarks.}
    We use clique--bridge--leaf motifs to isolate the loss of bridge-state information and chain-of-cliques benchmarks to study how local inter-community timing errors accumulate across several communities.
\end{enumerate}

Accordingly, the contribution of the paper is not simply to propose another multiscale epidemic model. Rather, it is to investigate a specific combination of microscopic CTMC simulation inside communities, macro-level stochastic transmission between communities, and explicit evaluation of when this separation of scales provides a useful approximation.

The central message of this paper is that local exactness and global accuracy are distinct properties. 
A two-scale simulator may preserve exact CTMC/Gillespie dynamics inside each community and still introduce structural bias at the interfaces between communities. 
This bias arises because the macro layer replaces boundary-node information, such as the infection status of bridge nodes, by aggregate community-level fractions. 
The resulting approximation error is therefore not a failure of the local Gillespie engines, but a consequence of the coarse-grained representation of inter-community transmission. 
We analyze this effect explicitly on clique--bridge--leaf motifs and then measure how the corresponding local timing errors accumulate in chain-of-communities benchmarks.



\section{Mathematical Model and Architecture of the Two-Scale Stochastic Simulator}

\subsection{Full-Network CTMC Reference Model}
\label{sec:full_network_ctmc_reference}

Before introducing the two-scale approximation, we define the exact microscopic process on the full network. This process is the reference model against which all later approximation errors are measured.

Let \(G=(V,E)\) be the full microscopic contact network. Each node has state in \(\{S,I,R\}\), and the node set is partitioned into communities
\[
V=\bigsqcup_{i=1}^{M} V_i .
\]
For a configuration \(x\in\{S,I,R\}^{V}\), the infection propensity along an oriented edge \((u,v)\in E\) is
\[
a_{u\to v}(x)
=\beta_{uv}\,\mathbf{1}\{x_u=I,\ x_v=S\}.
\]
The recovery propensity of node \(u\in V\) is
\[
a_u^{\mathrm{rec}}(x)
=\gamma_u\,\mathbf{1}\{x_u=I\}.
\]
Here \(\beta_{uv}\) is the edge-level transmission rate and \(\gamma_u\) is the node-level recovery rate. For a node \(v\), let \(x^{v:I}\) denote the configuration obtained from \(x\) by changing node \(v\) to state \(I\). Similarly, \(x^{u:R}\) denotes the configuration obtained by changing node \(u\) to state \(R\).

The full-network CTMC is the process with generator
\[
\begin{aligned}
(\mathcal L f)(x)
&=
\sum_{(u,v)\in E}
\beta_{uv}\,\mathbf{1}\{x_u=I,\ x_v=S\}
\bigl[f(x^{v:I})-f(x)\bigr]
\\
&\quad+
\sum_{u\in V}
\gamma_u\,\mathbf{1}\{x_u=I\}
\bigl[f(x^{u:R})-f(x)\bigr],
\end{aligned}
\]
for bounded test functions \(f\). This generator gives the exact full-network Gillespie/CTMC reference process. The two-scale simulator introduced below is therefore not defined in isolation: it is an approximation of this reference process, obtained by preserving local within-community CTMC dynamics while replacing the full collection of cross-community infection channels by macro-level import mechanisms.

\subsection{Two-Scale Simulator}
\label{sec:two_scale_simulator}

The two-scale simulator is defined on the fixed community partition
\[
V=\bigsqcup_{i=1}^{M}V_i,
\qquad N_i=|V_i|.
\]
For each community \(i\), let \(E_{ii}\) denote the within-community edges and let
\[
E_{ij}=\{(u,v)\in E:\ u\in V_i,\ v\in V_j\},
\qquad i\neq j,
\]
denote the oriented inter-community edges from community \(i\) to community \(j\) in the full reference network.

At time \(t\), the node-level state in community \(i\) is summarized by the sets
\[
S_i(t)=\{v\in V_i:X_v(t)=S\},\quad
I_i(t)=\{v\in V_i:X_v(t)=I\},\quad
R_i(t)=\{v\in V_i:X_v(t)=R\}.
\]
We distinguish these sets from their cardinalities,
\[
s_i(t)=|S_i(t)|,\qquad i_i(t)=|I_i(t)|,\qquad r_i(t)=|R_i(t)|,
\]
and from the normalized macro summaries
\[
\bar{s}_i(t)=\frac{s_i(t)}{N_i},\qquad
\bar{i}_i(t)=\frac{i_i(t)}{N_i},\qquad
\bar{r}_i(t)=\frac{r_i(t)}{N_i}.
\]
Thus \(S_i(t),I_i(t),R_i(t)\) are node sets, \(s_i(t),i_i(t),r_i(t)\) are counts, and \(\bar{s}_i(t),\bar{i}_i(t),\bar{r}_i(t)\) are fractions.

Each community has a local micro CTMC engine running on the induced graph \(G_i=(V_i,E_{ii})\). Between external imports, this engine simulates exactly the within-community infection and recovery channels. The macro layer observes only the community summaries, not the full boundary-node configuration. For each ordered pair \(i\neq j\), it assigns the effective inter-community hazard
\[
\lambda_{ij}(t)
=\beta_{\mathrm{out}}\,W_{ij}\,\bar{i}_i(t)\,\bar{s}_j(t),
\]
where \(W_{ij}\) is the aggregate coupling weight from community \(i\) to community \(j\), and \(\beta_{\mathrm{out}}\) is the macro-level transmission-rate parameter.

Let
\[
\Theta(t)=\sum_{i\neq j}\lambda_{ij}(t)
\]
be the total macro hazard. When a macro event occurs at time \(t\) and \(\Theta(t)>0\), the ordered source--destination pair is selected according to
\[
\mathbb P\{(i,j)\text{ is selected}\mid\text{macro event at }t\}
=\frac{\lambda_{ij}(t)}{\Theta(t)}.
\]
After selecting \(i\to j\), the simulator chooses a target node in the destination community using an import placement rule
\[
\Pi_{ij}(\cdot\mid t)
\quad\text{on}\quad S_j(t).
\]
For example, a simple rule is uniform placement over susceptible nodes,
\[
\Pi_{ij}(v\mid t)=\frac{1}{s_j(t)},\qquad v\in S_j(t),
\]
when \(s_j(t)>0\). More structured rules may use degree, boundary information, or an edge-conditioned oracle distribution. Once a target \(v\in S_j(t)\) is selected, the destination microstate is updated by the imported jump \(v:S\to I\), and the local propensities in community \(j\) are recomputed.

The synchronization mechanism couples these parts in continuous time. Micro engines advance their local CTMCs and export updated summaries \((\bar{i}_i(t),\bar{s}_i(t))\). The macro layer accumulates the total hazard \(\Theta(t)\) from these summaries. If the macro hazard budget is reached, the simulator interpolates the macro event time, selects the source--destination pair, applies the import placement rule, and updates the affected micro engine. This handshake is the mechanism by which the two-scale process replaces explicit cross-community CTMC channels by macro-generated imports.

\subsection{Micro Layer: Intra-Community CTMC Dynamics}
\ReadStatusTwo{}{Yulian}{1}{Michael}{0}

The micro simulator is the local stochastic engine of the two-scale framework. Its role is to represent intra-community dynamics(Poisson process) explicitly on the contact graph of a single community, while the interaction between communities is handled separately at the macro layer. In the present work, the micro layer follows the logic of classical continuous-time stochastic Gillespie algorithm~\cite{gillespie1977exact}, and is based on the efficient network-oriented implementation described in detail in~\cite{kuryliak2025simulating}.

Let community \(i\) be represented by the local graph \(G_i=(V_i,E_{ii})\), where nodes are individuals and edges are potential transmission contacts inside the community. At time \(t\), each node is in one of the states \(S\), \(I\), \(R\). The local state may therefore be written either as the node-level configuration \(X_i(t)\) or through the sets \(S_i(t)\), \(I_i(t)\), and \(R_i(t)\) defined above. The micro layer evolves through elementary stochastic events. For the internal SIR dynamics considered here, these are infection and recovery (transfers to other states is modeled with the same logic):
\[
S+I \xrightarrow{\beta_{\mathrm{in}}} I+I,
\qquad
I \xrightarrow{\gamma} R,
\]
where \(\beta_{\mathrm{in}}\) is the within-community infection rate and \(\gamma\) is the recovery rate.

The Gillespie construction is based on event rates. For each susceptible node \(v\in S_i(t)\), and each infectious node \(u\in I_i(t)\), let
\[
\lambda^{\mathrm{inf}}_v(t)=\beta_{\mathrm{in}}\, m_v(t),
\quad
\lambda^{\mathrm{rec}}_u(t)=\gamma.
\]
where \(m_v(t)\) is the number of infectious neighbors of \(v\) at time \(t\). Thus, the infection rate of a node is determined by its current local infectious neighborhood. Coeffisients \(\beta_{\mathrm{in}}\) and \(\gamma\) are the basic within-community infection and recovery rates.

Summing over nodes gives the type-specific local rates
\[
\theta_i^{\mathrm{inf}}(t)=\sum_{v\in S_i(t)} \lambda^{\mathrm{inf}}_v(t),
\qquad
\theta_i^{\mathrm{rec}}(t)=\sum_{u\in I_i(t)} \lambda^{\mathrm{rec}}_u(t)=\gamma i_i(t),
\]
Since all infectious nodes recover with the same rate \(\gamma\), this recovery selection is uniform over the currently infectious nodes. 

Hence the total local event rate
\[
\theta_i(t)=\theta_i^{\mathrm{inf}}(t)+\theta_i^{\mathrm{rec}}(t).
\]
Here \(\theta_i(t)\) is the total event rate of community \(i\) at the micro layer.

Given the current local state, the waiting time to the next internal event is sampled from the exponential law
\[
\Delta t_i=\frac{-\ln U}{\theta_i(t)},
\qquad U\sim\mathcal U(0,1).
\]
A second uniform random variable \(U_2\sim\mathcal U(0,1)\) is then used to choose the event type according to the relative propensities:
\[
\text{infection if } U_2\,\theta_i(t) < \theta_i^{\mathrm{inf}}(t),
\qquad
\text{recovery otherwise.}
\]

Once the event type has been selected, a specific node $e$ is chosen at random with probability proportional to its corresponding rate \(\lambda^{\mathrm{event}}_e(t)\). Thus, for an event of the given type,
\[
\mathbb P(\text{event on node } e \text{ next}\mid \mathrm{event})=
\frac{\lambda^{\mathrm{event}}_e(t)}{\theta_i^{\mathrm{event}}(t)}.
\]

After the event has been chosen, the local state is updated and only the affected node rates need to be refreshed. This event-driven update is the main advantage of the efficient Gillespie implementation used in \cite{kuryliak2025simulating}: rather than rebuilding the full set of propensities after each event, the simulator updates only those quantities changed by the local state transition. In this way, the micro layer preserves exact stochastic intra-community dynamics while remaining suitable as the local building block of the larger two-scale simulator.


\subsection{Macro Layer: Inter-Community Transmission}
\ReadStatusTwo{Macro Layer: Inter-Community Transmission}{Yulian}{0}{Michael}{0}

At the macro layer, the system is represented as a stochastic process on the graph of communities. Each community is treated as a meta-node, and the microscopic inter-community contacts of the original network are replaced by weighted macro-links that summarize the effective coupling between communities. In this way, the macro graph does not attempt to resolve individual cross-community transmissions one by one. Instead, it describes inter-community transmission through aggregate stochastic events between meta-nodes, while the detailed epidemic evolution inside each community remains delegated to the micro layer.

For every community \(i\), the micro layer exports the normalized state descriptors
\[
\bar{i}_i(t)=\frac{i_i(t)}{N_i},
\qquad
\bar{s}_i(t)=\frac{s_i(t)}{N_i}.
\]
Here \(\bar{i}_i(t)\) denotes the current contagiousness of community \(i\), measured through its infectious fraction, and \(\bar{s}_i(t)\) denotes its current susceptibility, measured through its susceptible fraction. If \(W_{ij}\) is the weight of the macro-link from community \(i\) to community \(j\), then the instantaneous rate at which community \(i\) seeds an infection in community \(j\) is modeled as
\[
\lambda_{ij}(t)
=\beta_{\mathrm{out}}\,W_{ij}\,\bar{i}_i(t)\,\bar{s}_j(t),
\]
where \(\beta_{\mathrm{out}}\) is the macro-level infection-rate parameter. If needed, this expression may be supplemented by a calibration factor to align macro-layer transmission with the corresponding micro-layer scale, but the central structure of the model is already captured by the formula above. 

The total macro intensity is
\[
\Theta(t)=\sum_{i\neq j}\lambda_{ij}(t).
\]
When a macro event fires, the ordered pair \(i\to j\) is selected with probability \(\lambda_{ij}(t)/\Theta(t)\). The destination node is then selected from \(S_j(t)\) according to the import placement distribution \(\Pi_{ij}(\cdot\mid t)\), and the selected susceptible node is changed to state \(I\). This placement rule is part of the simulator definition, because different choices of \(\Pi_{ij}\) can lead to different subsequent local CTMC paths.

\paragraph{Modeling assumption.}
The macro rate \(\lambda_{ij}(t)\) is treated as an effective inter-community transmission intensity. 
It is not assumed to be an exact coarse-graining of the full microscopic inter-community edge set unless this is stated explicitly. 
In particular, the formula assumes that the retained macro variables \(\bar{i}_i(t)\), \(\bar{s}_j(t)\), and \(W_{ij}\) contain enough information to approximate the aggregate transmission pressure from community \(i\) to community \(j\). 
When inter-community transmission depends on the states of specific boundary or bridge nodes, this assumption may fail and becomes a source of approximation error.

It is important to interpret this rate as an effective macro-level approximation rather than, in general, an exact coarse-graining of the full microscopic network. 
The exact inter-community intensity depends on which specific boundary nodes are infectious and which cross-community edges are active. 
By contrast, the macro rate retains only aggregate community-level quantities. 
This distinction is central for the accuracy analysis in Section~\ref{sec:accuracy_error}.

% Summing over all ordered pairs of distinct communities gives the total macro rate,
% \[
% \Theta(t)=\sum_{i\neq j}\lambda_{ij}(t).
% \]

A macro event has a simple interpretation: one imported infection is seeded in the destination community \(j\). Thus, the macro layer does not specify which exact individual transmitted infection across communities; rather, it represents the cumulative effect of inter-community coupling through community-to-community stochastic jumps. A crucial point is that the macro rates are not static. Even if no macro event occurs, the micro-layer states within communities continue to evolve, so the quantities \(\bar{i}_i(t)\) and \(\bar{s}_j(t)\), and hence the rates \(\lambda_{ij}(t)\), change continuously in time. Consequently, the macro-event process should not be described as an ordinary NHPP with deterministic intensity. Conditional on the micro-layer trajectories, and therefore on the predictable path of \(\Theta(t)\), macro events are generated by the cumulative-hazard construction. Unconditionally, the macro layer is a point process with stochastic predictable intensity \(\Theta(t)\). Because \(\Theta(t)\) varies with time, the waiting-time construction must be generalized beyond the classical fixed-rate Gillespie draw, which motivates the next subsection.

\subsection{Conditional NHPP Construction and Relation to Classical Gillespie}
\ReadStatusTwo{Hazard-Integral Timing and Relation to Classical Gillespie}{Yulian}{1}{Michael}{0}

The classical Gillespie algorithm generates the waiting time to the next event exactly when the total event rate is constant. If the total rate is \(\theta\), then the waiting time is exponentially distributed with parameter \(\theta\), and may be sampled as
\[
\Delta t = \frac{-\ln U}{\theta}, \qquad U\sim\mathcal U(0,1).
\]
This formula can also be interpreted through accumulated hazard. In the constant-rate case, the accumulated hazard grows linearly with time, so the event time is obtained by waiting until the linear hazard reaches a random threshold. In this sense, the usual Gillespie draw already has an integral interpretation, even though for constant rates the integral reduces to a simple division.

At the macro layer considered here, however, the direct fixed-rate draw is not sufficient. The quantities \(\bar{i}_i(t)\) and \(\bar{s}_j(t)\) continue to evolve because the micro layer changes the state of each community even between two macro events. As a consequence, the inter-community rates \(\lambda_{ij}(t)\), and therefore
\[
\Theta(t)=\sum_{i\neq j}\lambda_{ij}(t),
\]
are not constant. Conditional on the realized micro-layer trajectories, and therefore on the predictable path of \(\Theta(t)\), macro events are generated by the usual cumulative-hazard construction. Unconditionally, the macro layer is not an ordinary NHPP with deterministic intensity; it is a point process with stochastic predictable intensity \(\Theta(t)\). Thus, conditionally on the micro-layer paths, we measure elapsed risk through the \textbf{cumulative hazard}
\[
H(t)=\int_{t_0}^{t}\Theta(s)\,ds,
\]
where \(t_0\) is the time of the previous macro event. We then draw one uniform random variable \(U\sim\mathcal U(0,1)\) and define
\[
E=-\ln U,
\]
which we interpret as a random \textbf{hazard budget}.

The next macro event occurs at the first time \(t^\ast>t_0\) such that
\[
H(t^\ast)=E.
\]
Intuitively, the process keeps accumulating hazard over time, and the event fires exactly when this accumulated hazard reaches its random budget. This is the inverse, or integral, form of the usual Gillespie time-generation rule, generalized to dynamically varying rates. The classical formula is recovered as a sanity check: if \(\Theta(t)\equiv \theta\), then
\[
H(t)=\theta \cdot (t-t_0),
\]
and solving \(H(t^\ast)=E\) gives
\[
t^\ast-t_0=\frac{E}{\theta}=\frac{-\ln U}{\theta},
\]
which is exactly the standard Gillespie waiting-time expression. This formulation therefore preserves the logic of Gillespie's SSA while extending it to the stochastic, microstate-dependent macro rates required by the two-scale model.


\subsection{Synchronization Between Layers}
\ReadStatusTwo{Synchronization Between Layers}{Yulian}{1}{Michael}{0}

The two layers of the simulator are coupled through a bidirectional handshake. Synchronization is necessary because the macro layer does not possess an autonomous state of its own: its transmission intensities depend on the evolving epidemic state inside the communities, while the micro layer must react whenever the macro layer generates an imported infection. The purpose of synchronization is therefore to keep the aggregate inter-community process and the local intra-community CTMCs consistent in time.

In the forward direction, each community is first advanced at the micro layer over a synchronization interval \(dt\) (or, equivalently, over a micro interval \(\tau_{\mathrm{micro}}\)). At the end of that interval, the community exports refreshed summary quantities \(\bar{i}_i(t)\) and \(\bar{s}_i(t)\).

Thus, after each micro update, the macro layer receives a new snapshot of the current inter-community transmission pressure and continues hazard accumulation using the updated total rate. In the present formulation there is one current macro hazard budget \(E\), associated with the total hazard \(\Theta(t)\), rather than separate budgets for individual macro-links.

If the accumulated macro hazard does not reach the budget during the current synchronization interval, the simulation simply advances to the next synchronization point, where the micro states are refreshed again and the procedure repeats. If, however, the cumulative hazard reaches the budget inside the current interval, then the macro event occurs at an interpolated time \(t^\ast\) within that interval. At that moment, one inter-community transmission is realized, a destination community is selected according to the current macro intensities, and one imported infection is inserted into the corresponding micro model. Conceptually, this means that the current interval is split: the destination community evolves up to \(t^\ast\), receives the external infection at \(t^\ast\), and then continues from the updated state over the remainder of the interval.

This backward coupling requires a consistent update of the local micro-layer clock. Suppose that, before the import, the destination community was evolving from local time \(t_0\) with total internal rate \(\theta_i(t_0)\) and a local random budget
\[
E_i=-\ln U_i, \qquad U_i\sim \mathcal U(0,1).
\]
Under the usual Gillespie logic, the next internal event would have been scheduled at
\[
t_i^{\mathrm{old}}=t_0+\frac{E_i}{\theta_i(t_0)},
\]
because the local state, and hence the rate, remain constant between two consecutive local events. If an imported infection arrives earlier, at time \(t^\ast<t_i^{\mathrm{old}}\), then the previously scheduled local event is interrupted. Up to time \(t^\ast\), the local process has already consumed the hazard
\[
H_i(t^\ast)=\int_{t_0}^{t^\ast}\theta_i(s)\,ds.
\]
Since \(\theta_i(s)=\theta_i(t_0)\) on \([t_0,t^\ast)\) in the standard Gillespie setting, this becomes
\[
H_i(t^\ast)=\theta_i(t_0)\,(t^\ast-t_0).
\]
The remaining local budget after the import is therefore
\[
E_i^{\mathrm{res}}=E_i-\theta_i(t_0)\,(t^\ast-t_0).
\]

At time \(t^\ast\), the imported infection changes the local state instantaneously, so the old rate is no longer valid. Let \(\theta_i^{\mathrm{new}}(t^\ast)\) denote the updated total local rate after the import has been applied and the local propensities have been recomputed. The new next-event time, denoted by \(t_i^{\mathrm{new}}\), is then obtained by continuing from the remaining budget under the updated rate:
\[
t_i^{\mathrm{new}}=t^\ast+\frac{E_i^{\mathrm{res}}}{\theta_i^{\mathrm{new}}(t^\ast)}.
\]
% Equivalently, in integral form, \(t_i^{\mathrm{new}}\) is the first time satisfying
% \[
% \int_{t^\ast}^{t_i^{\mathrm{new}}}\theta_i^{\mathrm{new}}(s)\,ds = E_i^{\mathrm{res}}.
% \]
Thus, \(t_i^{\mathrm{new}}\) is not drawn from scratch as if the elapsed time before \(t^\ast\) had never happened; rather, it is computed by subtracting the already consumed hazard from the original budget and then propagating the local clock forward under the updated state. If another external import arrives before \(t_i^{\mathrm{new}}\), the same residual-budget update is applied again. In this way, the synchronization rule is not only an exchange of summary statistics, but also a time-consistency condition: the macro layer advances by accumulated hazard built from the evolving micro summaries, while each imported macro event instantaneously modifies the destination CTMC and forces a consistent recalculation of its local next-event time.


\section{--------------In Progress---------------}

\section{Accuracy and Approximation Error of the Two-Scale Simulator}
\label{sec:accuracy_error}

\subsection{Accuracy Framework and Evaluation Logic}
\ReadStatusTwo{Accuracy Framework and Evaluation Logic}{Yulian}{1}{Michael}{0}

The key point of the accuracy analysis is that the two-scale approximation does not primarily modify the local stochastic dynamics inside communities. 
Instead, it modifies the interface between communities. 
Thus, the central question is not whether the local Gillespie engines are exact, but whether the macro layer contains enough information to reproduce the timing and placement of inter-community events generated by the full-network CTMC.

The purpose of this part of the paper is to assess when the two-scale simulator provides an accurate representation of the underlying continuous-time stochastic dynamics on the original graph, and to identify the mechanisms that govern its approximation error. Throughout this analysis, the full-network CTMC simulated by the Gillespie algorithm is treated as the reference process. The proposed simulator is therefore evaluated as an approximation to this microscopic reference: it preserves explicit intra-community dynamics at the micro layer, while replacing the full set of inter-community infection channels by a macro-level abstraction driven by community-level state summaries.

In this paper, \emph{accuracy} is understood relative to the reference CTMC over a fixed time horizon. The comparison focuses on stochastic observables that are relevant for community-structured dynamics. These include inter-community transmission times, the ordering of arrivals across communities, prevalence trajectories within communities and at the network level, and the variability of these quantities across sample paths. Thus, accuracy is not reduced to agreement of mean trajectories alone. A useful abstraction should also reproduce, within the resolution of its macro state, the uncertainty and event-level stochastic patterns generated by the reference process.

This evaluation logic is motivated by the role of rare and heterogeneous inter-community transmissions in community-structured graphs. Two simulators may produce similar mean prevalence curves while differing in the timing or conditional structure of inter-community infections. Such differences can change the subsequent local trajectories inside affected communities and may therefore alter the distribution of later observables. For this reason, the present work does not restrict attention to deterministic limits or mean-field summaries. Instead, it asks whether the two-scale abstraction captures the relevant stochastic behavior of the full CTMC at the level of event times, community-level paths, and derived quantities such as arrival times and outbreak sizes.

The remainder of this part is organized around benchmark cases designed to expose specific aspects of the approximation. We begin with analytically tractable graph motifs where the inter-community mechanism can be inspected directly at the edge level. We then move to structured community-based graph families in which the interaction between intra-community dynamics and inter-community transmission is shaped by graph topology. These benchmarks provide a controlled setting in which discrepancies between the two-scale simulator and the full-network reference can be attributed to identifiable structural features.

\paragraph{Standing assumptions for the accuracy analysis.}
Unless stated otherwise, the full-network process is a continuous-time Markov chain on a fixed finite graph \(G=(V,E)\), with node states in \(\{S,I,R\}\) or in \(\{S,I\}\) for the SI benchmark cases. 
The node set is partitioned into fixed communities
\[
V=\bigsqcup_{i=1}^{M}V_i.
\]
Within-community transitions are generated by the usual Gillespie/SSA propensities on the induced subgraphs \(G_i=G[V_i]\). 
Inter-community edges are represented in the full reference process by the sets \(E_{ij}\), while the two-scale simulator replaces these edge-level channels by macro-level import events. 
All stochastic clocks are exponential conditional on the current state, and all rates are assumed constant between consecutive CTMC events unless a time-dependent macro hazard is explicitly specified. 
The full-network Gillespie CTMC is used as the reference process for all accuracy statements.

\subsection{Sources of Approximation Error}
\ReadStatusTwo{Sources of Approximation Error}{Yulian}{1}{Michael}{0}

The two-scale simulator is designed to preserve the local CTMC dynamics within each community while replacing the original inter-community mechanism by a macro-level abstraction and a coupling rule. At the modeling level, it departs from the full-network Gillespie process through the way inter-community transmissions are represented and injected into the micro layer. At the implementation level, additional numerical error may arise from synchronization and hazard evaluation.

Let
\[
V=\bigsqcup_{i=1}^{M}V_i
\]
be the fixed partition of the original graph into communities. Let \(G_i=G[V_i]\) be the subgraph induced by community \(i\), and let
\[
E_{ij}=\{(u,v)\in E:\ u\in V_i,\ v\in V_j\}
\]
denote the set of inter-community edges from community \(i\) to community \(j\). Denote by \(X(t)\) the full-network CTMC state and by \(X_i(t)\) its restriction to \(V_i\). In the full-network process, communities are coupled through edge-level inter-community events in the sets \(E_{ij}\). In the two-scale simulator, the same coupling is represented by a macro-generated import process into the local community models.

\paragraph{Proposition (Conditional micro-layer exactness).}
Fix the initial microstate in each community. Also fix the complete import process into every community, including all import times and the target nodes in the destination communities. Conditional on this import process, the micro layer in the two-scale simulator evolves each \(X_i(t)\) exactly as the Gillespie CTMC on \(G_i\), with the same intra-community transitions and with the specified imports applied as exogenous \(S\to I\) jumps. Consequently, any approximation error relative to the full-network Gillespie reference arises from how the import process is generated and coupled, not from the local CTMC simulation itself.

\paragraph{Proof.}
Once the import times and target nodes are fixed, the evolution inside each community is determined by the same local transition rules as in the full-network CTMC restricted to \(G_i\). Between two exogenous imports, the micro layer applies the Gillespie algorithm to the same intra-community reaction channels and propensities. At an import time, the prescribed target node is updated by the same exogenous state transition. Therefore the sample-path law of \(X_i(t)\), conditional on the fixed import process, coincides with the corresponding local CTMC law. The remaining discrepancy can only come from the construction of the import process itself, including its rate, target placement, and any numerical treatment of its timing.

\paragraph{Interpretation.}
This proposition formalizes the main structural message of the paper: within-community exactness does not imply exactness of the full two-scale process. 
The local Gillespie engines are exact conditional on the imports they receive, but the imports themselves are generated by a coarse-grained macro mechanism. 
Therefore, the dominant approximation error is an interface error between the micro and macro layers. In what follows, we decompose this discrepancy into three sources: the rate abstraction that replaces the microscopic inter-community intensity, the import placement rule that replaces the edge-conditioned target choice, and the numerical timing error introduced by synchronization and hazard evaluation.

\subsubsection{Rate Abstraction Error}
\ReadStatusTwo{Rate Abstraction Error}{Yulian}{1}{Michael}{0}

In the full-network CTMC, the instantaneous inter-community infection intensity from community \(i\) to community \(j\) is determined by the active cross-community infection opportunities in the original graph. For a full microscopic state \(x\), the exact intensity is
\begin{equation}
\label{eq:true_inter_intensity_error}
a^\star_{ij}(x)
=
\beta_{\mathrm{out}}\sum_{(u,v)\in E_{ij}}
w_{uv}\,\mathbf{1}\{x_u=I,\ x_v=S\}.
\end{equation}
The two-scale simulator replaces this edge-level quantity by the macro-rate
\begin{equation}
\label{eq:macro_inter_intensity_error}
\lambda_{ij}(x)
=
\beta_{\mathrm{out}}\, W_{ij}\,\bar{i}_i(x)\,\bar{s}_j(x),
\qquad
\bar{i}_i(x)=\frac{i_i(x)}{N_i},
\quad
\bar{s}_j(x)=\frac{s_j(x)}{N_j},
\end{equation}
where \(W_{ij}\) is an aggregate weight associated with the meta-edge \((i,j)\). The macro state therefore retains only community-level fractions, not the exact boundary configuration of the original graph.

\paragraph{Proposition (Rate abstraction).}
Define the rate abstraction error process by
\[
\varepsilon_{ij}(t)
=
a^\star_{ij}(X(t))-\lambda_{ij}(X(t)).
\]
If the macro-rate \(\lambda_{ij}(x)\) is exact for all full-network states in a given state space, then the microscopic cross-community intensity \(a^\star_{ij}(x)\) must be completely determined by the retained aggregates \((\bar{i}_i(x),\bar{s}_j(x),W_{ij})\). Equivalently, if there exist two states \(x\) and \(y\) with
\[
\bar{i}_i(x)=\bar{i}_i(y),\qquad
\bar{s}_j(x)=\bar{s}_j(y),
\]
and the same fixed \(W_{ij}\), but with different boundary or bridge configurations such that
\[
a^\star_{ij}(x)\neq a^\star_{ij}(y),
\]
then no macro-rate of the form \eqref{eq:macro_inter_intensity_error} can be exact for both states.

\paragraph{Proof.}
The macro-rate \(\lambda_{ij}(x)\) is a function only of \(\bar{i}_i(x)\), \(\bar{s}_j(x)\), \(W_{ij}\), and the parameter \(\beta_{\mathrm{out}}\). Therefore, two full states with identical values of these quantities necessarily receive the same macro-rate. If their microscopic intensities \(a^\star_{ij}\) differ, then the same value of \(\lambda_{ij}\) cannot equal both microscopic intensities. Hence exactness of the macro-rate over the relevant state space requires the true inter-community intensity to be determined by the retained macro variables.

\paragraph{Interpretation.}
Equation~\eqref{eq:true_inter_intensity_error} depends on which specific nodes in \(V_i\) and \(V_j\) are incident to inter-community edges and on their current states. The macro summary \((\bar{i}_i,\bar{s}_j)\) discards this bridge-level structure. Thus, the approximation is reliable only when the boundary nodes are sufficiently representative of the communities with respect to inter-community exposure, or when the graph structure makes the aggregate summaries informative enough for the true inter-community intensity. When inter-community edges are concentrated on a small subset of boundary nodes, or when infection status is strongly correlated with boundary position, the error process \(\varepsilon_{ij}(t)\) can be substantial.

\subsubsection{Import Placement Error}
\ReadStatusTwo{Import Placement Error}{Yulian}{1}{Michael}{0}

In the full-network Gillespie process, an inter-community event \(i\to j\) occurs through a specific active edge. Conditional on observing such an event at state \(x\), the target node in \(V_j\) is not chosen only from the community-level susceptible fraction. It is determined by the firing edge \((u,v)\in E_{ij}\). Provided \(a^\star_{ij}(x)>0\), the full CTMC induces the edge-conditioned target distribution
\begin{equation}
\label{eq:exact_import_placement}
\pi^\star_{ij}(v\mid x)
=
\frac{
\mathbf{1}\{x_v=S\}
\sum_{u\in V_i:\ (u,v)\in E_{ij}}
\beta_{uv}\,\mathbf{1}\{x_u=I\}
}
{a^\star_{ij}(x)},
\qquad v\in V_j.
\end{equation}
Here \(\beta_{uv}=\beta_{\mathrm{out}}w_{uv}\) under the weighted-edge parametrization in \eqref{eq:true_inter_intensity_error}.
Consequently, the induced distribution over target nodes \(v\in V_j\) depends on the inter-community adjacency of \(v\) and on the states of its cross-community neighbors.

The two-scale simulator instead injects an imported infection into \(V_j\) by an abstracted community-level seeding rule. Let \(\Pi_{ij}(t)\) denote the simulator's conditional distribution over susceptible target nodes in \(V_j\), given that a macro event \(i\to j\) fires at time \(t\). A simple choice is uniform seeding over susceptible nodes in \(V_j\), but the following criterion applies to any abstracted seeding distribution.

\paragraph{Proposition (Import placement).}
Fix a state \(x\) and condition on the event that an inter-community transmission \(i\to j\) occurs at that state. The import placement mechanism in the two-scale simulator reproduces the immediate post-event microstate law of the full-network CTMC if and only if the seeding distribution \(\Pi_{ij}(\cdot\mid x)\) coincides with the full-network conditional distribution \(\pi^\star_{ij}(\cdot\mid x)\) in \eqref{eq:exact_import_placement}. If these distributions differ, then the two-scale simulator changes the distribution of the destination microstate immediately after import.

\paragraph{Proof.}
Conditioning on the occurrence of an \(i\to j\) event fixes the event type at the community level but not the target node in \(V_j\). In the full-network CTMC, the post-event state in \(V_j\) is obtained by infecting the node selected by the edge-level firing distribution \(\pi^\star_{ij}(\cdot\mid x)\). In the two-scale simulator, it is obtained by infecting a node sampled from \(\Pi_{ij}(\cdot\mid x)\). The two post-event laws coincide exactly when these two target-node distributions coincide. Otherwise, at least one susceptible node receives different probability under the two mechanisms, and the law of the post-event microstate is perturbed.

\paragraph{Interpretation.}
Import placement error is a coupling error, not a minor implementation detail. Even if the macro event time were correct, seeding the infection into a uniformly random susceptible node rather than into the edge-conditioned target can modify the early intra-community dynamics after import. This effect is expected to be most visible in motifs where inter-community connectivity is concentrated on a small set of boundary nodes, because the full-network process necessarily targets that subset whereas an abstracted rule may distribute imports more broadly.

\subsubsection{Synchronization and Numerical Approximation Error}
\ReadStatusTwo{Synchronization and Numerical Approximation Error}{Yulian}{1}{Michael}{0}

The structural errors above arise from the two-scale abstraction itself. A separate, secondary source of error may arise from the numerical realization of synchronization. The macro layer generates inter-community events from the total rate
\[
\Theta(t)
=
\sum_{i\neq j}
\lambda_{ij}(t).
\]
In the ideal continuous-time formulation, if \(\Theta(t)\) were known exactly and the macro event time \(t^\star\) were obtained from
\[
\int_{t_0}^{t^\star}
\Theta(s)\,ds
=
E,
\qquad
E\sim \mathrm{Exp}(1),
\]
then the event time would be exact conditional on the realized micro-layer trajectories that determine \(\Theta(t)\). Unconditionally, this is a macro point process with stochastic predictable intensity \(\Theta(t)\), not a deterministic-intensity NHPP. In practice, however, \(\Theta(t)\) is updated through synchronization between layers, and the cumulative hazard is accumulated numerically.

\paragraph{Proposition (Synchronization timing error bound).}
Let
\[
H(t)
=
\int_{t_0}^{t}
\Theta(s)\,ds
\]
be the exact cumulative hazard, and let \(\widehat{H}(t)\) be the numerically accumulated hazard used by the implementation. Let \(t^\star\) and \(\widehat{t}\) be defined by
\[
H(t^\star)=E,
\qquad
\widehat{H}(\widehat{t})=E.
\]
Assume that, on an interval containing \(t^\star\) and \(\widehat{t}\),
\[
|H(t)-\widehat{H}(t)|\le \eta
\]
and
\[
\Theta(t)\ge \theta_{\min}>0.
\]
Then the event-time error is bounded by
\[
|\widehat{t}-t^\star|
\le
\frac{\eta}{\theta_{\min}}.
\]

\paragraph{Proof.}
Since \(H\) is increasing with derivative \(\Theta(t)\ge\theta_{\min}\), we have
\[
|H(\widehat{t})-H(t^\star)|
\ge
\theta_{\min}|\widehat{t}-t^\star|.
\]
On the other hand, \(H(t^\star)=E=\widehat{H}(\widehat{t})\), so
\[
|H(\widehat{t})-H(t^\star)|
=
|H(\widehat{t})-\widehat{H}(\widehat{t})|
\le \eta.
\]
Combining the two inequalities gives the bound.

\paragraph{Interpretation.}
The resulting timing discrepancy is controlled by the error in approximating the cumulative hazard \(H\), including quadrature over synchronization intervals, interpolation of the threshold crossing time, and any delay in updating \(\bar{i}_i(t)\) and \(\bar{s}_j(t)\) relative to the underlying micro-layer evolution. Synchronization and numerical timing error is not the primary approximation mechanism of the two-scale simulator. It can, however, add timing error on top of the structural discrepancies described above. For this reason, the benchmark analysis focuses primarily on rate abstraction and import placement, while the numerical implementation of hazard accumulation must remain consistent with the intended stochastic-intensity macro point process.


\subsection{Clique--bridge--leaf benchmark}
\ReadStatusTwo{Clique--bridge--leaf benchmark}{Yulian}{1}{Michael}{0}
The clique--bridge--leaf benchmark is used as a minimal counterexample to the idea that aggregate community fractions are always sufficient for inter-community transmission. 
Although the source community is internally well mixed as a clique, the export event depends on the infection status of a small bridge set. 
This makes the motif analytically simple while still exposing the central interface error of the two-scale approximation.

This benchmark is introduced as a minimal test for identifying exactly what is lost when passing from the full microscopic network to the two-scale abstraction. In the previous sections, the micro layer preserves the exact CTMC/Gillespie dynamics inside each community, while the macro layer describes inter-community transmission through the aggregated intensity
\[
\lambda_{ij}(t)=\beta_{\mathrm{out}} W_{ij}\bar{i}_i(t)\bar{s}_j(t).
\]
It is precisely this difference between the exact states of individual nodes and the aggregated community-level variables that is central for the present benchmark.

\subsubsection{Graph construction and motivation}
\ReadStatusTwo{Graph construction and motivation}{Yulian}{1}{Michael}{0}

Consider two communities. Community \(0\) is the infection source and has local graph
\[
G_0=(V_0,E_0).
\]
In this benchmark, it is a clique
\[
K_n,
\]
so that
\[
N_0=n
\]
nodes belong to the source community. Community \(1\) is the receiver and consists of a single leaf node \(A\). This node has no internal dynamics before the import event; it can only become infected through an inter-community contact from community \(0\).

For analytical clarity, this benchmark uses the irreversible SI version of the local process, so that once a bridge node becomes infected, it remains infectious.

Inside the clique \(K_n\), we distinguish a bridge set
\[
B\subset K_n,
\qquad |B|=k.
\]
The leaf \(A\) is connected only to the nodes in \(B\). These bridge nodes are therefore the only microscopic nodes of community \(0\) through which infection can reach \(A\). In the full microscopic network, transmission from community \(0\) to community \(1\) is possible not from any infected node of the clique, but only from infected bridge nodes.

Let
\[
B_I(t)
=
\sum_{b\in B}
\mathbf{1}\{X_b(t)=I\}
\]
be the number of infected bridge nodes at time \(t\). This is the key microscopic quantity for the export event: the leaf is infected by the concrete states of the bridge nodes, not by the average infected fraction in the clique.

\subsubsection{Exact Micro dynamics}
\ReadStatusTwo{Exact Micro dynamics}{Yulian}{1}{Michael}{0}

\paragraph{Proposition. Exact bridge-dependent micro hazard.}
In the full Micro model, the instantaneous import intensity into the leaf node \(A\) is
\[
a_{01}^{\mathrm{Micro}}(t)
=
\beta_{\mathrm{out}} B_I(t),
\]
provided that node \(A\) is still susceptible. If \(A\) is already infected, further import into it is not considered.

This statement follows directly from the local graph structure. The leaf node \(A\) has contacts only with bridge nodes. Each infected bridge node creates one active infectious edge to \(A\) with rate \(\beta_{\mathrm{out}}\). If there are \(B_I(t)\) infected bridge nodes, then the total intensity of these parallel edges is \(\beta_{\mathrm{out}}B_I(t)\).

This formula has an important interpretation. Even if many nodes in the clique are already infected, if no bridge node has yet been infected, then
\[
B_I(t)=0
\]
and therefore
\[
a_{01}^{\mathrm{Micro}}(t)=0.
\]
Thus, in the exact microscopic model, the leaf node is fully protected until the infection actually reaches at least one bridge node.

Let the time of first bridge activation be denoted by
\[
T_B=\inf\{t:B_I(t)>0\}.
\]
In this benchmark, it is also convenient to interpret it as the bridge-finding time:
\[
T_{\mathrm{find}}=T_B.
\]

\subsubsection{Two-phase structure of the exact process}
\ReadStatusTwo{Two-phase structure of the exact process}{Yulian}{1}{Michael}{0}

\paragraph{Proposition. Two-phase structure for a single bridge and a single leaf.}
Assume \(k=1\), SI dynamics, and no recovery. Let \(T_B\) be the infection time of the bridge node and let \(T_A\) be the infection time of the leaf node. Then the exact Micro process has the decomposition
\[
T_A=T_B+Z,
\qquad
Z\sim \mathrm{Exp}(\beta_{\mathrm{out}}),
\]
where \(Z\) is the post-bridge waiting time, conditionally on the history up to \(T_B\). Consequently,
\[
\mathbb{E}[T_A]
=
\mathbb{E}[T_B]
+
\frac{1}{\beta_{\mathrm{out}}}.
\]

The first phase is the delay phase. The infection spreads inside the clique, but the leaf node still has no active infectious contact. Over the whole interval
\[
0\leq t<T_B,
\]
the risk to the leaf node is zero.

The second phase begins after the unique bridge node enters state \(I\). From that moment, the leaf node sees exactly one active bridge--leaf edge with intensity \(\beta_{\mathrm{out}}\). Since in the SI model the infected bridge node does not recover, the subsequent dynamics inside the clique no longer change the hazard for this one leaf node. Therefore, after the bridge node becomes infected, the additional waiting time is exponential with mean
\[
\frac{1}{\beta_{\mathrm{out}}}.
\]

This decomposition highlights an important property of the exact Micro model: the size of the clique and its internal dynamics influence the infection time of the leaf node only through the delay phase, that is, through the time needed to infect the bridge node. Once the unique bridge node is already infected, the remaining events inside the clique no longer affect the expected infection time of this one leaf node; it is determined only by a single active edge with rate \(\beta_{\mathrm{out}}\).

For several bridge nodes, this interpretation remains useful, but the second phase may already change. If the leaf node has contacts with several bridge nodes, then after additional bridge nodes become infected, the intensity becomes
\[
a_{01}^{\mathrm{Micro}}(t)=\beta_{\mathrm{out}}B_I(t).
\]
That is, several infected bridge nodes create several parallel active edges to the leaf node. In that case, the post-bridge time is no longer necessarily equal to \(1/\beta_{\mathrm{out}}\), because the later infection of additional bridge nodes may increase the hazard and shorten the second phase.

Let the import time into the leaf node \(A\) be denoted by
\[
T_A.
\]
Then the exact infection clock of the leaf node is given through the accumulated hazard:
\[
\int_0^{T_A}
\beta_{\mathrm{out}}B_I(s)\,ds
=
-\ln U,
\qquad
U\sim\mathcal{U}(0,1).
\]
Up to time \(T_B\), this integral does not accumulate, because \(B_I(s)=0\). After the first bridge node becomes infected, it starts growing. If further bridge nodes become infected later, the accumulation speed increases, since \(B_I(t)\) becomes larger.

\subsubsection{Bridge-finding time in a clique}
\ReadStatusTwo{Bridge-finding time in a clique}{Yulian}{1}{Michael}{0}

\paragraph{Proposition. Expected bridge-finding time in a clique.}
For the case of one bridge node, a simple analytical check can be obtained. If \(k=1\), then the expected time until the bridge node becomes infected is
\[
\mathbb{E}[T_B]
=
\frac{H_{n-1}}{\beta_{\mathrm{in}} n},
\]
where
\[
H_{n-1}
=
\sum_{m=1}^{n-1}\frac{1}{m}
\]
is the harmonic number.

After the bridge node becomes infected, the leaf node receives one active infectious contact with rate \(\beta_{\mathrm{out}}\), so the additional expected time until import is
\[
\frac{1}{\beta_{\mathrm{out}}}.
\]
Hence,
\[
\mathbb{E}[T_A^{\mathrm{Micro}}]
=
\frac{H_{n-1}}{\beta_{\mathrm{in}} n}
+
\frac{1}{\beta_{\mathrm{out}}}.
\]

For an arbitrary number of bridge nodes \(k\), it is convenient to consider the time of first activation of the bridge set. Its expected value can be written as a sum over clique states:
\[
\mathbb{E}[T_B]
=
\sum_{m=1}^{n-k}
\frac{\binom{n-k}{m}}{\binom{n}{m}}
\cdot
\frac{1}{\beta_{\mathrm{in}} m(n-m)}.
\]
Here the factor
\[
\frac{\binom{n-k}{m}}{\binom{n}{m}}
\]
is the probability that among the first \(m\) infected nodes of the clique there is still no bridge node. If \(k\) increases, this probability decreases rapidly, so the first bridge node is activated earlier. In the limiting case of large \(k\), the exact hazard becomes closer to the macro hazard, because the bridge nodes better represent the whole community.

\subsection{Mean-field versus exact hazard bias}
\ReadStatusTwo{Mean-field versus exact hazard bias}{Yulian}{1}{Michael}{0}

This subsection explains why an apparently reasonable macro hazard may still distort event timing. 
The macro hazard can be correct as a conditional mean of the exact microscopic bridge hazard, but event times depend on accumulated hazards and survival probabilities, not only on instantaneous expected rates. 
Consequently, replacing the random bridge-dependent hazard by its mean-field counterpart can change the distribution of export times.

The clique--bridge--leaf benchmark separates the exact microscopic mechanism from its aggregate approximation. The exact model depends on the state of the bridge nodes, while the macro model sees only the community-level infected fraction \(\bar{i}_0(t)\). This subsection makes this difference explicit and explains why replacing the exact hazard by a mean-field hazard changes the timing of export events.

\subsubsection{Macro approximation of bridge risk}
\ReadStatusTwo{Macro approximation of bridge risk}{Yulian}{1}{Michael}{0}

For comparison with the exact Micro process, the macro layer does not see the specific bridge nodes. It receives only the aggregated infected fraction \(\bar{i}_0(t)\). If we interpret the inter-community weight \(W_{01}\) as corresponding to the number of potential bridge contacts, that is,
\[
W_{01}=k,
\]
then the macro rate for import from community \(0\) to community \(1\) takes the form
\[
\lambda_{01}^{\mathrm{Macro}}(t)
=
\beta_{\mathrm{out}} k \bar{i}_0(t)\bar{s}_1(t).
\]
Since in the benchmark community \(1\) consists of a single susceptible leaf node before import, we may take
\[
\bar{s}_1(t)=1.
\]
Therefore, the macro rate simplifies to
\[
\lambda_{01}^{\mathrm{Macro}}(t)
=
\beta_{\mathrm{out}} k\frac{i_0(t)}{n}.
\]
This is the aggregated version of the bridge risk. It replaces the exact number of infected bridge nodes \(B_I(t)\) by its mean-field estimate:
\[
B_I(t)
\approx
k\frac{i_0(t)}{n}.
\]

The corresponding macro clock is given by
\[
\int_0^{T_A^{\mathrm{Macro}}}
\beta_{\mathrm{out}} k\frac{i_0(s)}{n}\,ds
=
-\ln U.
\]

\subsubsection{Conditional-mean interpretation}
\ReadStatusTwo{Conditional-mean interpretation}{Yulian}{1}{Michael}{0}

\paragraph{Proposition. Macro hazard is the conditional mean of the exact micro hazard.}
At the level of the instantaneous mean, the macro rate is a natural approximation of the exact micro hazard:
\[
\mathbb{E}
\left[
a_{01}^{\mathrm{Micro}}(t)
\mid i_0(t)
\right]
=
\lambda_{01}^{\mathrm{Macro}}(t).
\]

Indeed, if we only know the number of infected nodes in the clique \(i_0(t)\), but not their identities, then the probability that a particular bridge node is infected equals
\[
\frac{i_0(t)}{n}.
\]
Since there are \(k\) bridge nodes, the expected number of infected bridge nodes is
\[
\mathbb{E}[B_I(t)\mid i_0(t)]
=
k\frac{i_0(t)}{n}.
\]
Then
\[
\mathbb{E}
\left[
a_{01}^{\mathrm{Micro}}(t)
\mid i_0(t)
\right]
=
\mathbb{E}
\left[
\beta_{\mathrm{out}} B_I(t)
\mid i_0(t)
\right]
=
\beta_{\mathrm{out}} k\frac{i_0(t)}{n}
=
\lambda_{01}^{\mathrm{Macro}}(t).
\]

Thus, the macro rate is the correct conditional mean of the exact micro hazard. However, this does not mean that the macro model reproduces the exact infection-time distribution. The reason is that the import time is determined not by the hazard itself, but by the exponential of the accumulated hazard.

\subsubsection{Bias in event timing}
\ReadStatusTwo{Bias in event timing}{Yulian}{1}{Michael}{0}

The early-export bias should be formulated at the level of accumulated hazards, not only at the level of instantaneous hazards. Let
\[
X_t
=
\int_0^t \beta_{\mathrm{out}}B_I(s)\,ds
\]
be the exact accumulated bridge hazard up to time \(t\). Let \(\mathcal F_t^I\) denote the information retained by the macro layer up to time \(t\), for example the aggregate infection-count path \(\{i_0(s):0\le s\le t\}\). Suppose that the macro accumulated hazard is the conditional mean
\[
\bar X_t
=
\mathbb{E}\!\left[X_t\mid \mathcal F_t^I\right].
\]
Under the exchangeability assumptions used above, this corresponds to
\[
\bar X_t
=
\int_0^t \beta_{\mathrm{out}}k\frac{i_0(s)}{n}\,ds.
\]

\paragraph{Proposition. Conditional Jensen survival inequality.}
The exact conditional survival probability satisfies
\[
\mathbb{E}\!\left[\exp(-X_t)\mid\mathcal F_t^I\right]
\geq
\exp(-\bar X_t).
\]
Consequently, after taking expectations,
\[
S^{\mathrm{Micro}}(t)
\geq
S^{\mathrm{Macro}}(t),
\]
where
\[
S^{\mathrm{Micro}}(t)
=
\mathbb{E}\!\left[\exp(-X_t)\right],
\qquad
S^{\mathrm{Macro}}(t)
=
\mathbb{E}\!\left[\exp(-\bar X_t)\right].
\]

\paragraph{Proof.}
The function \(x\mapsto \exp(-x)\) is convex. Applying Jensen's inequality conditionally on \(\mathcal F_t^I\) gives
\[
\mathbb{E}\!\left[\exp(-X_t)\mid\mathcal F_t^I\right]
\geq
\exp\!\left(
-\mathbb{E}[X_t\mid\mathcal F_t^I]
\right)
=
\exp(-\bar X_t).
\]
Taking expectations gives the unconditional survival inequality.

\paragraph{Corollary. Expected import-time ordering.}
If the survival inequality holds for all \(t\ge 0\) and the expectations are finite, then
\[
\mathbb{E}[T_A^{\mathrm{Micro}}]
\geq
\mathbb{E}[T_A^{\mathrm{Macro}}].
\]

\paragraph{Proof.}
Since the expectation of a nonnegative infection time can be represented as the area under the survival curve,
\[
\mathbb{E}[T_A]
=
\int_0^\infty S(t)\,dt,
\]
integrating the survival inequality over time gives the result.

Intuitively, this means that the macro model may create premature leakage of risk. As soon as infected nodes appear in the clique, the quantity \(i_0(t)/n\) becomes positive, and the macro layer begins accumulating hazard for import. But in the full microscopic network, import may still be impossible if the infected nodes are not bridge nodes.

The clique--bridge--leaf benchmark is useful precisely because it isolates structural abstraction error from other sources of error. There is no complicated community structure here, no many-path interaction, and no heterogeneous internal topology. There is only a clique, \(k\) bridge nodes, and one leaf receiver. If even in such a simple system the macro approximation deviates from the full Gillespie reference, then the reason for the deviation is clear: the macro layer does not know the exact states of the bridge nodes.

\subsection{Numerical benchmark on the clique--bridge--leaf graph}
\ReadStatusTwo{Numerical benchmark on the clique--bridge--leaf graph}{Yulian}{1}{Michael}{0}

\subsubsection{Experimental setup}
\ReadStatusTwo{Experimental setup}{Yulian}{1}{Michael}{0}

In the implemented numerical benchmark, two models were compared: the full Micro simulation on the explicit network and the MicroMacro approximation. The parameterization corresponds to the case \(n=100\), \(k=2\), that is, the source community is a clique of 100 nodes, and inter-community transmission passes through two bridge nodes.

Since the source community is a clique, community \(0\) reaches near-complete infection very quickly in both models. This is expected, because the number of active \(S-I\) edges is large during the initial and intermediate phases of the intra-community process. However, the main quantity of interest in this benchmark is not the visual shape of the infection curve inside the clique, but the timing of bridge activation and export. Therefore, the most informative comparison is given by the timing distributions in Figure~\ref{fig:clique_leaf_density_panels}.

\subsubsection{Timing distributions}
\ReadStatusTwo{Timing distributions}{Yulian}{1}{Michael}{0}

\begin{figure}[h]
    \centering
    \includegraphics[width=\linewidth]{Images/Metric_density_panels.png}
    \caption{Density-style distributions of bridge and export timings.}
    \label{fig:clique_leaf_density_panels}
\end{figure}

Figure~\ref{fig:clique_leaf_density_panels} separately shows the bridge infection time, the time from bridge infection to export, and the total time from simulation start to export. The first panel shows that the distributions of bridge infection times in Micro and MicroMacro are nearly identical. This is an important sanity check: both models reproduce a similar time scale for the activation of the bridge mechanism in the clique.

The second panel shows the distribution of the delay from bridge infection to export. It has a long right tail, corresponding to the exponential character of waiting for a rare inter-community transmission. The third panel shows the full time from the start to export; here too, the Micro and MicroMacro distributions are close.

\subsubsection{Implications for the two-scale approximation}
\ReadStatusTwo{Implications for the two-scale approximation}{Yulian}{1}{Michael}{0}

This result is important for the interpretation of the benchmark. Theoretically, the mean-field macro hazard may produce premature leakage of risk. However, in this specific configuration with \(n=100\) and \(k=2\), the effect is moderate. The reason is that the bridge nodes in the clique are activated rather early, while the main variability of the total export time is associated not with the bridge infection time, but with the long random delay between bridge activation and the export event itself. Thus, the difference between the exact and aggregate hazards is partly masked by the broader distribution of inter-community waiting times.

This does not negate the presence of structural bias. On the contrary, the benchmark shows in which regimes it may be small. If there are enough bridge nodes, if the clique mixes rapidly, and if the export rate is relatively small, then the macro approximation may agree well with the full Micro reference. In such a regime, the quantity
\[
B_I(t)
\]
quickly becomes close to its mean-field analogue
\[
k\frac{I_0(t)}{n}.
\]
Therefore, the macro hazard
\[
\lambda_{01}^{\mathrm{Macro}}(t)
=
\beta k\frac{I_0(t)}{n}
\]
becomes an acceptable effective description of the exact micro hazard.

However, the same benchmark also indicates the direction for stress tests. To reveal premature leakage more strongly, one should consider cases with smaller \(k\), slower internal spreading inside the community, or a higher inter-community transmission rate. In such regimes, the macro layer may begin accumulating risk substantially earlier than the full-network Gillespie process allows real export through the bridge nodes.

Thus, the clique--bridge--leaf benchmark serves two functions. First, it provides a simple analytical model for checking the consistency of the Micro and MicroMacro simulators. Second, it demonstrates that the macro rate based on \(\bar{i}_i(t)\) is not an exact coarse-graining, but an effective approximation. Its accuracy depends on how well the bridge nodes represent the state of the whole community. For well-mixed communities and sufficiently broad bridge contacts, the approximation may be practically accurate. For narrow boundary bottlenecks, it may systematically accelerate inter-community transmission.




\subsection{Chain--of--cliques benchmark}
\ReadStatusTwo{Chain--of--cliques benchmark}{Yulian}{0}{Michael}{0}



After the clique--bridge--leaf benchmark, the next natural test case is a chain of cliques. The previous benchmark analyzed a single local export event: the process had to leave one clique through a bridge node and activate one leaf node. In two-scale modeling, however, it is also important to understand how such local timing errors behave when the process propagates through several communities in sequence.

We consider a chain of \(L\) communities,
\[
G_0 \to G_1 \to \cdots \to G_{L-1}.
\]
Each community \(G_i=(V_i,E_i)\) is a clique of size
\[
N_i=n.
\]
The process starts in the first community \(G_0\). A downstream community \(G_{i+1}\) becomes active only after a successful export event from \(G_i\). Neighboring cliques are connected through a narrow interface, for example through \(k\) bridge nodes or through an equivalent set of sparse inter-community contacts.

This benchmark tests not an isolated local transition, but a propagation wave across a sequence of communities. If the timing error of one transition is small but repeated several times, it may accumulate along the chain. Therefore, the chain--of--cliques benchmark is a stronger test of the global timing error of the two-scale approximation.

\subsubsection{Arrival times and propagation delays}
\ReadStatusTwo{Arrival times and propagation delays}{Yulian}{0}{Michael}{0}

Let the number of active nodes in community \(i\) at time \(t\) be denoted by
\[
I_i(t).
\]
The first-arrival time in community \(i\) is defined as
\[
T_i=\inf\{t:I_i(t)>0\}.
\]
For the initially active community we set
\[
T_0=0.
\]

The local propagation delay between two neighboring communities is
\[
\Delta T_i=T_{i+1}-T_i,
\qquad i=0,\ldots,L-2.
\]
The global first-arrival time in the last community is
\[
T_{\mathrm{last}}=T_{L-1}.
\]

\paragraph{Proposition. Arrival time in the last community decomposes into local delays.}
The first-arrival time in the last community can be written as
\[
T_{\mathrm{last}}
=
\sum_{i=0}^{L-2}\Delta T_i.
\]
Consequently,
\[
\mathbb{E}[T_{\mathrm{last}}]
=
\sum_{i=0}^{L-2}
\mathbb{E}[\Delta T_i].
\]

This follows directly from the definition of the local propagation delays. Each \(\Delta T_i\) measures the time needed to move from community \(G_i\) to community \(G_{i+1}\). Summing all these delays gives the total time needed to reach \(G_{L-1}\).

If all transitions are parametrized identically, that is, all cliques have the same size \(n\), the same number of bridge nodes \(k\), the same transmission intensity \(\beta\), and each newly activated community starts from one active node, then all \(\Delta T_i\) have the same distribution. In this homogeneous case,
\[
\mathbb{E}[T_{\mathrm{last}}]
=
(L-1)\mathbb{E}[\Delta T].
\]
This identity is the main reason why the chain benchmark is useful: a local timing discrepancy of a single transition can scale with the number of transitions in the chain.

\subsubsection{Expected time of one exact Micro transition}
\ReadStatusTwo{Expected time of one exact Micro transition}{Yulian}{0}{Michael}{0}

A single transition
\[
G_i \to G_{i+1}
\]
can be decomposed into two phases. The first phase is the time needed for the process inside \(G_i\) to reach at least one bridge node. We denote this time by
\[
T_{\mathrm{find}}.
\]
The second phase is the time between the first bridge activation and the successful export event into the next community. We denote it by
\[
T_{\mathrm{cross}}.
\]

\paragraph{Proposition. One exact Micro transition decomposes into a bridge-finding phase and a crossing phase.}
The local propagation delay satisfies
\[
\Delta T
=
T_{\mathrm{find}}+T_{\mathrm{cross}},
\]
and therefore
\[
\mathbb{E}[\Delta T]
=
\mathbb{E}[T_{\mathrm{find}}]
+
\mathbb{E}[T_{\mathrm{cross}}].
\]

This decomposition separates two different mechanisms. The first term measures how long it takes the process to reach the inter-community interface. The second term measures how long it takes to cross this interface once at least one bridge node is active.

Ц
\paragraph{Proposition. Exact transition time for a single bridge node.}
If \(k=1\), then the expected time needed to find the bridge node in a clique is
\[
\mathbb{E}[T_{\mathrm{find}}]
=
\frac{H_{n-1}}{\beta n},
\]
where
\[
H_{n-1}
=
\sum_{m=1}^{n-1}\frac{1}{m}
\]
is the harmonic number.

After the unique bridge node becomes active, the next community sees one active inter-community edge with rate \(\beta\). Hence,
\[
\mathbb{E}[T_{\mathrm{cross}}]
=
\frac{1}{\beta}.
\]
Therefore, for one bridge node, the expected time of one exact Micro transition is
\[
\mathbb{E}[\Delta T]
=
\frac{H_{n-1}}{\beta n}
+
\frac{1}{\beta}.
\]

For a homogeneous chain, this gives the baseline expression
\[
\mathbb{E}[T_{\mathrm{last}}]
=
(L-1)
\left(
\frac{H_{n-1}}{\beta n}
+
\frac{1}{\beta}
\right).
\]
This formula shows that the expected arrival time in the last community is composed of identical local blocks. In each clique, the process first has to find the bridge node and then cross one active inter-community edge.

\subsubsection{Several bridge nodes}
\ReadStatusTwo{Several bridge nodes}{Yulian}{0}{Michael}{0}

For \(k>1\), the first phase has the same interpretation, but the bridge set contains several nodes.

\paragraph{Proposition. Expected bridge-finding time for a bridge set.}
Under the uniform seeding assumption inside each newly activated clique, the expected time until the first activation of the bridge set is
\[
\mathbb{E}[T_{\mathrm{find}}]
=
\sum_{m=1}^{n-k}
\frac{\binom{n-k}{m}}{\binom{n}{m}}
\cdot
\frac{1}{\beta m(n-m)}.
\]

Here \(m\) is the number of active nodes in the clique. The factor
\[
\frac{\binom{n-k}{m}}{\binom{n}{m}}
\]
is the probability that among the first \(m\) active nodes there is still no bridge node, while
\[
\frac{1}{\beta m(n-m)}
\]
is the expected residence time of the clique in the state with \(m\) active nodes.

The second phase is no longer identical to the single-bridge case. After the first bridge node becomes active, the export hazard is initially \(\beta\). If additional bridge nodes become active before the export event occurs, the export hazard increases:
\[
a_{\mathrm{export}}^{\mathrm{Micro}}(t)
=
\beta B_I(t),
\]
where \(B_I(t)\) is the number of active bridge nodes in the current community.

\paragraph{Proposition. Crossing time for several bridge nodes is bounded by the single-edge and fully active interface cases.}
For \(k>1\), the expected crossing time satisfies
\[
\frac{1}{\beta k}
\leq
\mathbb{E}[T_{\mathrm{cross}}]
\leq
\frac{1}{\beta}.
\]

The lower bound corresponds to the case where all \(k\) bridge nodes are already active and the export hazard is \(\beta k\). The upper bound corresponds to the case where only the first bridge node contributes to the export clock.

Hence, for \(k>1\), the expected time of one transition lies in the interval
\[
\mathbb{E}[T_{\mathrm{find}}]
+
\frac{1}{\beta k}
\leq
\mathbb{E}[\Delta T]
\leq
\mathbb{E}[T_{\mathrm{find}}]
+
\frac{1}{\beta}.
\]
Consequently, the expected arrival time in the last community satisfies
\[
(L-1)
\left(
\mathbb{E}[T_{\mathrm{find}}]
+
\frac{1}{\beta k}
\right)
\leq
\mathbb{E}[T_{\mathrm{last}}]
\leq
(L-1)
\left(
\mathbb{E}[T_{\mathrm{find}}]
+
\frac{1}{\beta}
\right).
\]
This analytical interval is sufficient for the main benchmark interpretation. It shows how the parameters \(L\), \(n\), \(k\), and \(\beta\) affect the expected propagation time without introducing a full finite-state recursion.

\subsubsection{Relation to the MicroMacro approximation}
\ReadStatusTwo{Relation to the MicroMacro approximation}{Yulian}{0}{Michael}{0}

In the MicroMacro model, the identities of the bridge nodes are not tracked explicitly. Instead, the macro layer observes only the aggregated contagiousness
\[
\bar{i}_i(t)=\frac{i_i(t)}{N_i}.
\]
For the transition
\[
G_i\to G_{i+1},
\]
the macro hazard has the form
\[
\lambda_{i,i+1}^{\mathrm{Macro}}(t)
=
\beta_{\mathrm{out}} W_{i,i+1}\bar{i}_i(t)\bar{s}_{i+1}(t).
\]
Before the next community is activated, its susceptible fraction is approximately
\[
\bar{s}_{i+1}(t)\approx 1.
\]
If the macro-link weight corresponds to the number of bridge contacts,
\[
w_{i,i+1}=k,
\]
then
\[
\lambda_{i,i+1}^{\mathrm{Macro}}(t)
=
\beta k\frac{I_i(t)}{n}.
\]
Thus, the MicroMacro model replaces the exact bridge-dependent hazard
\[
\beta B_I(t)
\]
by the aggregated mean-field estimate
\[
\beta k\frac{I_i(t)}{n}.
\]

\paragraph{Proposition. A local timing shift accumulates along a homogeneous chain.}
Suppose that, for each transition in a homogeneous chain, the MicroMacro approximation induces the same expected timing shift
\[
\delta
=
\mathbb{E}[\Delta T^{\mathrm{MicroMacro}}]
-
\mathbb{E}[\Delta T^{\mathrm{Micro}}].
\]
Then the expected shift in the arrival time of the last community is
\[
\mathbb{E}[T_{\mathrm{last}}^{\mathrm{MicroMacro}}]
-
\mathbb{E}[T_{\mathrm{last}}^{\mathrm{Micro}}]
=
(L-1)\delta.
\]

This proposition follows directly from the decomposition of \(T_{\mathrm{last}}\) into local delays. It explains why a chain benchmark is more sensitive than a single-transition benchmark: even a small local discrepancy can become visible after several consecutive transitions.

The clique--bridge--leaf benchmark showed that the mean-field replacement may produce premature leakage: the macro hazard can start accumulating before the exact Micro process has reached the bridge nodes. In a chain of cliques, the same local mechanism is repeated at several consecutive transitions. At the same time, the empirical timing shift of an implemented MicroMacro simulator may also reflect layer synchronization, macro-rate calibration, the import placement rule, and the numerical approximation of the hazard integral. For this reason, the sign and magnitude of the observed shift are reported through arrival-time distributions in the numerical benchmark.

\subsubsection{Representative trajectories}
\ReadStatusTwo{Representative trajectories}{Yulian}{0}{Michael}{0}

\textbf{The question how to choose representative curve is still open. Is median value at each point correct? On the plots we representative (solid) curve was seleted with the lowest kind of mean squared error.}

The available representative trajectories show the expected wave-like propagation pattern. After a community is activated, the local curve grows rapidly, which is consistent with clique dynamics. Thus, the main difference between Micro and MicroMacro is not the shape of the local growth curve inside a clique, but the timing of first activation:
\[
T_1,T_2,\ldots,T_{L-1}.
\]

\begin{figure}[t]
    \centering
    \includegraphics[width=\linewidth]{plots/k_5_n_100_inter_1_macro_chain_micro_complete_p0p031/Micro_vs_MicroMacro_metric_panels.png}
    \caption{Representative per-community trajectories for the Micro and MicroMacro models on the chain--of--cliques benchmark.}
    \label{fig:chain_cliques_community}
\end{figure}

Figure~\ref{fig:chain_cliques_community} compares the community-level trajectories of the two models. Both models activate communities in the same sequential order, and each activated clique rapidly approaches its saturated state. The visible difference is primarily a shift in arrival times. In the shown representative run, several MicroMacro community curves are shifted to the right relative to the corresponding Micro curves, indicating a delayed arrival in those communities.

\begin{figure}[t]
    \centering
    \includegraphics[width=\linewidth]{plots/k_5_n_100_inter_1_macro_chain_micro_complete_p0p031/Micro_vs_MicroMacro.png}
    \caption{Aggregate comparison of Micro and MicroMacro trajectories with interquartile bands.}
    \label{fig:chain_cliques_aggregate}
\end{figure}

Figure~\ref{fig:chain_cliques_aggregate} gives an aggregate comparison with interquartile bands. The bands indicate that a single representative trajectory is not sufficient to determine the systematic direction of the approximation error. The relevant quantities are the empirical distributions of the arrival times and propagation delays.

For each random variable
\[
X\in\{\Delta T_i,T_{\mathrm{last}}\},
\]
we compare the empirical distribution functions
\[
F_X^{\mathrm{Micro}}(t),
\qquad
F_X^{\mathrm{MicroMacro}}(t).
\]
Useful distributional metrics include the Wasserstein-1 distance
\[
W_1(X)=
\int_0^\infty
\left|
F_X^{\mathrm{Micro}}(t)
-
F_X^{\mathrm{MicroMacro}}(t)
\right|\,dt,
\]
the Kolmogorov--Smirnov distance
\[
D_{\mathrm{KS}}(X)=
\sup_t
\left|
F_X^{\mathrm{Micro}}(t)
-
F_X^{\mathrm{MicroMacro}}(t)
\right|,
\]
the mean shift
\[
\Delta \mu_X
=
\mu_X^{\mathrm{MicroMacro}}
-
\mu_X^{\mathrm{Micro}},
\]
and the quantile shift
\[
\Delta q_p(X)
=
q_p^{\mathrm{MicroMacro}}(X)
-
q_p^{\mathrm{Micro}}(X).
\]
The median shift \(p=0.5\) measures the typical timing displacement, while \(p=0.9\) is useful for detecting changes in the late-arrival tail.

If
\[
\Delta \mu_{T_{\mathrm{last}}}<0,
\]
then the MicroMacro approximation reaches the last community earlier on average. If
\[
\Delta \mu_{T_{\mathrm{last}}}>0,
\]
then it reaches the last community later on average. This sign is reported from the empirical statistics of the benchmark rather than from a single trajectory.

\subsubsection{Monte Carlo arrival-time comparison}
\ReadStatusTwo{Monte Carlo arrival-time comparison}{Yulian}{1}{Michael}{0}

The numerical comparison was run as a paired Monte Carlo experiment. Each replication uses the same network seed and stochastic run seed for the Micro and MicroMacro simulators, but the metrics below are distributional rather than single-trajectory errors. The experiment uses a chain of complete communities with
\[
L=5,\qquad n=100,\qquad k=1,\qquad
\beta=0.005,\qquad \gamma=0,
\]
\(\beta_{\mathrm{macro}}=\beta\), \(\tau_{\mathrm{micro}}=1\), macro multiplier \(T=1\), terminal time \(T_{\mathrm{end}}=1000\), and initial infection at node \(0\) in community \(0\). We ran \(100\) paired replications. The run seeds were \(42,\ldots,141\), and the network seeds were \(42,\ldots,141\), since the network was regenerated for each replication.

For each run and model, the output records
\[
T_0,T_1,\ldots,T_4,\qquad
\Delta T_i=T_{i+1}-T_i,\quad i=0,\ldots,3,
\qquad
T_{\mathrm{last}}=T_4.
\]
The raw replication metadata, per-replication arrival table, wide arrival table, summary metrics, comparison metrics, paired error table, and runtime table were written to
\[
\texttt{data/chain\_cliques\_arrival\_comparison/chain\_L5\_n100\_k1\_beta0p005\_paper\_100reps/}.
\]
The generated figures were written to
\[
\texttt{plots/k\_5\_n\_100\_inter\_1\_macro\_chain\_micro\_complete\_p0p031/}.
\]

\begin{table}[t]
    \centering
    \begin{tabular}{lrrrrr}
        \hline
        Observable & \(W_1\) & \(D_{\mathrm{KS}}\) & \(\Delta\mu\) & \(\Delta q_{0.5}\) & \(\Delta q_{0.9}\) \\
        \hline
        \(T_{\mathrm{last}}\) & 48.750 & 0.115 & -23.811 & -32.752 & -64.636 \\
        \(\Delta T_0\) & 38.490 & 0.137 & 36.362 & 24.475 & 154.750 \\
        \(\Delta T_1\) & 24.466 & 0.107 & -7.062 & 10.267 & -80.546 \\
        \(\Delta T_2\) & 29.127 & 0.124 & -16.604 & -22.270 & -14.657 \\
        \(\Delta T_3\) & 37.446 & 0.218 & -29.647 & -62.726 & 20.791 \\
        \hline
    \end{tabular}
    \caption{Distributional comparison of Micro and MicroMacro arrival observables in the chain--of--cliques experiment. Shifts are MicroMacro minus Micro. Metrics are computed on completed observations; completion probabilities are reported in the text.}
    \label{tab:chain_cliques_arrival_metrics}
\end{table}

By \(T_{\mathrm{end}}=1000\), the last community was reached in \(90\) of \(100\) Micro runs and in \(64\) of \(100\) MicroMacro runs. Among completed terminal arrivals, the Micro mean was \(676.431\), the MicroMacro mean was \(652.620\), and the mean shift was therefore \(-23.811\). The corresponding medians were \(655.979\) and \(623.227\), giving a median shift of \(-32.752\). The \(0.9\)-quantiles were \(985.588\) for Micro and \(920.952\) for MicroMacro. Thus, conditional on completion before the time horizon, MicroMacro has slightly earlier terminal arrivals; however, its lower completion probability is itself an important discrepancy and indicates a heavier censored late-arrival component in this finite-horizon experiment.

\begin{figure}[t]
    \centering
    \includegraphics[width=0.85\linewidth]{plots/k_5_n_100_inter_1_macro_chain_micro_complete_p0p031/chain_cliques_tlast_cdf_paper_100reps.png}
    \caption{Empirical CDF of completed \(T_{\mathrm{last}}\) observations for the Micro and MicroMacro models in the 100-replication chain--of--cliques experiment.}
    \label{fig:chain_cliques_tlast_cdf}
\end{figure}

\begin{figure}[t]
    \centering
    \includegraphics[width=\linewidth]{plots/k_5_n_100_inter_1_macro_chain_micro_complete_p0p031/chain_cliques_selected_delay_density_paper_100reps.png}
    \caption{Density-style histograms of selected local propagation delays \(\Delta T_i\). The panels show the first, middle, and final transition in the chain.}
    \label{fig:chain_cliques_delay_cdfs}
\end{figure}

\begin{figure}[t]
    \centering
    \includegraphics[width=0.95\linewidth]{plots/k_5_n_100_inter_1_macro_chain_micro_complete_p0p031/chain_cliques_error_growth_paper_100reps.png}
    \caption{Layer-wise growth of arrival-time discrepancy. The panels report \(W_1(T_i)\), mean shift \(\Delta\mu(T_i)\), and quantile shifts \(\Delta q_{0.5}(T_i)\), \(\Delta q_{0.9}(T_i)\) as functions of community index.}
    \label{fig:chain_cliques_error_growth}
\end{figure}

The layer-wise metrics show that the discrepancy is not a constant local offset repeated at every transition. The Wasserstein distance for \(T_i\) is \(38.490\) at community \(1\), \(44.188\) at community \(2\), \(55.945\) at community \(3\), and \(48.750\) at the terminal community. The mean shift changes sign: MicroMacro is delayed at \(T_1\) and \(T_2\), with shifts \(36.362\) and \(23.156\), but the terminal completed-arrival shift is \(-23.811\). The local delay metrics also have mixed signs: the mean shift is \(36.362\) for \(\Delta T_0\), \(-7.062\) for \(\Delta T_1\), \(-16.604\) for \(\Delta T_2\), and \(-29.647\) for \(\Delta T_3\). Therefore the discrepancy grows through the middle of the chain and then remains of the same order, while its signed bias changes direction.

The measured simulation runtime was \(28.765\) seconds for the \(100\) Micro runs and \(96.519\) seconds for the \(100\) MicroMacro runs. With the speedup convention
\[
\mathrm{Speedup}
=
\frac{
\mathrm{runtime}_{\mathrm{Micro}}
}{
\mathrm{runtime}_{\mathrm{MicroMacro}}
},
\]
this gives \(\mathrm{Speedup}=0.298\). Thus, in this small benchmark and current Python implementation, MicroMacro was slower than the full Micro simulator. This runtime result should not be interpreted as a general complexity statement for large multiscale systems; it reflects the present implementation, the small graph size, and the cost of synchronizing the macro hazard at \(\tau_{\mathrm{micro}}=1\).

This experiment identifies the net arrival-time discrepancy between the implemented Micro and MicroMacro simulators, but it does not uniquely attribute the discrepancy to one mechanism. The mean-field bridge abstraction can create premature macro hazard accumulation from the aggregate infected fraction before the exact bridge node is infected. At the same time, the implemented comparison also includes layer synchronization, the choice \(\beta_{\mathrm{macro}}=\beta\), and the import placement rule in the destination clique. The lower MicroMacro completion probability by \(T_{\mathrm{end}}\) also shows that the finite time horizon matters. Therefore the result should be read as a full-system benchmark. A direct ablation should repeat the same experiment while varying \(\tau_{\mathrm{micro}}\), recalibrating \(\beta_{\mathrm{macro}}\), increasing \(T_{\mathrm{end}}\) or applying an explicit right-censoring estimator, and replacing uniform import placement by bridge-aware or edge-conditioned placement where the code supports it.


------------------------------------------------------------------------------------------------------------------------------------------------------------------------------------------------------------------------------------------------------------------------------------------------------------------------------------------------------------------------------------------------
\section{Other notes, ideas, generated sections:}
------------------------------------------------------------------------------------------------------------------------------------------------------------------------------------------------------------------------------------------------------------------------------------------------------------------------------------------------------------------------------------------------


\section{Benchmark Case I: Clique-to-Leaf Transmission}
\label{sec:benchmark_clique_leaf}

\subsection{Graph Construction and Motivation}
\ReadStatusTwo{Graph Construction and Motivation}{Yulian}{1}{Michael}{0}

The first benchmark is designed to isolate the simplest nontrivial export mechanism from a dense community. We consider a complete graph \(K_N\), interpreted as the source community, and select one node \(b\in K_N\) as the bridge node. A single leaf node \(\ell\) is attached to this bridge node by one external edge \((b,\ell)\). In the numerical experiment considered below, the source community is a clique with \(N=100\) nodes and the target consists of one leaf.

This motif separates two different parts of the process. The first part is the time until the bridge node becomes infected inside the dense source community. The second part is the time from bridge infection to export through the single bridge-to-leaf edge. This decomposition is useful because the first part depends on intra-community dynamics, whereas the second part depends only on the active external edge once the bridge node is infected and the leaf is still susceptible. Thus, the clique-to-leaf benchmark provides a transparent test of whether the two-scale simulator reproduces the relevant event-level behavior of the full-network Gillespie reference.

We denote by \(T_b\) the infection time of the bridge node and by \(T_{\ell}\) the infection time of the leaf. The three quantities used in the comparison are therefore \(T_b\), the export delay \(T_{\ell}-T_b\), and the total export time \(T_{\ell}\). These observables make it possible to distinguish an error in the internal infection of the bridge node from an error in the export mechanism itself.

\subsection{Theoretical Export-Time Argument}
\ReadStatusTwo{Theoretical Export-Time Argument}{Yulian}{1}{Michael}{0}

The key analytical observation is that, after the bridge node has become infectious, the waiting time to infect the leaf is controlled by the single active bridge-to-leaf edge. Let \(b_{\mathrm{out}}\) denote the transmission rate along this edge. In the simplest case, \(b_{\mathrm{out}}\) is the edge-level transmission rate; more generally, it may include the corresponding edge weight.

\paragraph{Proposition (Export delay after bridge infection).}
Assume that at time \(T_b\) the bridge node \(b\) is infectious, the leaf \(\ell\) is susceptible, and the bridge-to-leaf edge remains active until the export event is realized. Then
\[
T_{\ell}-T_b \sim \mathrm{Exp}(b_{\mathrm{out}}),
\]
and hence
\[
\mathbb{E}\!\left[T_{\ell}-T_b\right]=\frac{1}{b_{\mathrm{out}}}.
\]
In particular, under these assumptions, the distribution of the export delay does not depend on how many additional infections occur inside the clique after time \(T_b\).

\paragraph{Proof.}
Condition on the bridge infection time \(T_b\). After this time, the leaf can be infected only through the edge \((b,\ell)\). Under the stated assumptions, this edge contributes a constant hazard \(b_{\mathrm{out}}\) for the transition of the leaf from susceptible to infected. Internal infection events inside the clique may change the state of other clique nodes, but they do not change the hazard of the single edge \((b,\ell)\), provided that \(b\) remains infectious and \(\ell\) remains susceptible. Therefore the waiting time from \(T_b\) to the leaf infection is exponential with rate \(b_{\mathrm{out}}\), which gives the stated expectation.

\paragraph{Interpretation.}
This proposition is not a general exactness theorem for arbitrary community-structured graphs. It applies to the simplified clique-to-leaf motif under the stated conditions. Its role is to identify a quantity that can be tested directly: once the bridge node has been infected, the export delay should be governed by the bridge-to-leaf hazard rather than by the subsequent number of infections inside the clique.

\subsection{Numerical Validation of Bridge and Export Timing}
\ReadStatusTwo{Numerical Validation of Bridge and Export Timing}{Yulian}{1}{Michael}{0}

We compare the full microscopic Gillespie simulator with the two-scale simulator on the clique-to-leaf benchmark. The comparison is distributional rather than only mean-based. We examine the bridge infection time \(T_b\), the post-bridge export delay \(T_{\ell}-T_b\), and the total export time \(T_{\ell}\). These three quantities correspond to the three stages of the benchmark: internal activation of the bridge node, export from the bridge to the leaf, and the resulting total arrival time.

\begin{figure}[t]
    \centering
    \includegraphics[width=\textwidth]{Images/Metric_density_panels.png}
    \caption{
    Density-style distributions of bridge and export timings in the clique-to-leaf benchmark. 
    The source community is a clique of \(N=100\) nodes connected to a single leaf through a bridge node. 
    The left panel shows the bridge infection time, the middle panel shows the delay from bridge infection to export, and the right panel shows the total time from initialization to export. 
    The comparison between the full microscopic Gillespie simulation and the two-scale simulator evaluates whether the abstraction preserves the relevant export-time behavior in this analytically transparent motif.
    }
    \label{fig:clique_leaf_export_density}
\end{figure}

Figure~\ref{fig:clique_leaf_export_density} shows that the empirical density curves of the microscopic and two-scale simulations are close for all three observables. The left panel indicates that the distribution of bridge infection times is reproduced at the level relevant for this benchmark. The middle panel isolates the export delay after bridge infection, which is the direct numerical counterpart of the theoretical argument above. The right panel combines these two components and compares the total time from initialization to export.

The comparison should be interpreted as evidence of agreement for this controlled benchmark, not as a general proof of exactness. In particular, the figure supports the claim that the two-scale abstraction can reproduce the bridge-to-leaf export mechanism when the source community is dense and the external target is structurally simple. A table of summary statistics, such as empirical means, standard deviations, and relative errors for the three timing observables, can be added later to complement the density plots.

\subsection{Interpretation of the Approximation Error}
\ReadStatusTwo{Interpretation of the Approximation Error}{Yulian}{1}{Michael}{0}

The clique-to-leaf benchmark is favorable to the two-scale abstraction for two reasons. First, the source community is dense, so aggregate quantities such as the infected fraction are informative about the internal progression of the process. Second, the export mechanism is concentrated on one structurally explicit edge. This makes it possible to separate the approximation of the bridge infection time from the approximation of the export delay.

In terms of the error mechanisms introduced earlier, the middle panel of Figure~\ref{fig:clique_leaf_export_density} mainly tests whether the export hazard after bridge infection is represented correctly. If the bridge-to-leaf rate is captured by the macro representation, then the post-bridge export delay should agree with the full-network reference. The right panel then tests the composition of the two components: the internal time to infect the bridge and the external time from bridge infection to export.

The import placement error is limited in this particular benchmark because the target consists of a single leaf. There is no nontrivial distribution over multiple destination nodes to approximate. Similarly, synchronization error should remain secondary provided that hazard accumulation and event-time interpolation are implemented consistently. Thus, this benchmark serves as a baseline case: it verifies the export mechanism in a dense and structurally simple setting before moving to more complex multi-community benchmarks, such as chains of cliques and chains of connected random graphs. 






\section{The Baseline of Transmission: The Single Edge and the Leaf}

\paragraph{The Atomic Unit of a Network Epidemic} 
Before addressing the complexities of densely connected communities and macroscopic abstractions, it is essential to build our mathematical intuition from the ground up. To do this, we must examine the most fundamental atomic unit of any network-based epidemic model: a single, isolated pair of individuals. 

Imagine a network consisting of exactly two nodes. The first node acts as the source and is currently in the \textit{Infected} state. The second node, which we will call the "leaf," is in the \textit{Susceptible} state. These two nodes are connected by a single, directed edge that represents the physical contact or pathway through which the disease can spread.



\paragraph{The Mechanics of Transmission Time}
In continuous-time stochastic models, such as the Gillespie algorithm, an infection does not jump across an edge instantaneously or at deterministic intervals. Instead, the transmission is modeled as a Poisson process. The edge possesses a specific transmission rate, denoted as $\beta$, which represents the expected number of successful transmission events per unit of time (e.g., per day). 

Because the epidemic process is memoryless—meaning the probability of infection in the next minute does not depend on how long the leaf has already been exposed—the actual time $T$ it takes for the infection to successfully cross the edge and infect the leaf is a continuous random variable. Specifically, $T$ strictly follows an exponential distribution governed by the rate parameter $\beta$. 

\paragraph{The Expected Time to Infection}
The defining characteristic of an exponentially distributed random variable is its mean. If the transmission rate is $\beta$, mathematical statistics dictate that the expected time (or average waiting time) until the Susceptible leaf becomes Infected is precisely the inverse of the rate:
$$ \mathbb{E}[T] = \frac{1}{\beta} $$
For instance, if the transmission rate $\beta$ is $0.1$ per day, we intuitively expect the infection to successfully cross the edge in $1/0.1 = 10$ days on average. This simple expected value, $\mathbb{E}[T] = 1/\beta$, is the absolute baseline truth of the Susceptible-Infected (SI) model.

\paragraph{Modeling the Process: The Hazard Integral}
How do stochastic simulators computationally determine the exact moment $T$ when the leaf becomes Infected? They do so by tracking the continuous accumulation of "infection pressure" or risk. 

In the simulation, the leaf node remains Susceptible as long as the cumulative risk it absorbs from its Infected neighbor does not exceed a randomly generated biological threshold. This threshold is drawn from a probability distribution by taking $-\ln(U)$, where $U$ is a uniformly distributed random variable such that $U \sim \text{Uniform}(0, 1)$. 

The time of infection $T$ is mathematically defined as the exact moment when the integral of the transmission rate (the hazard rate) over time equals this stochastic threshold:
$$ \int_{0}^{T} \beta \, dt = -\ln(U) $$
Because the source node is permanently Infected and $\beta$ is constant, this integral simplifies easily to $\beta \cdot T = -\ln(U)$, which rearranging gives $T = -\ln(U)/\beta$. This is the standard method for computationally sampling a random transmission time from an exponential distribution.

\paragraph{The Macro-Scale Objective}
Understanding this simple atomic interaction is crucial for the rest of our analysis. The expected time $\mathbb{E}[T] = 1/\beta$ represents the fundamental physical reality of a single connection. 

When we transition to a two-scale simulator, we replace the single Infected source node with an entire complex community containing $n$ nodes, and we abstract the links into macroscopic edges. Our ultimate mathematical goal—and the core challenge of this paper—is to ensure that the hazard integral designed for these macroscopic community-to-community edges mathematically preserves this exact baseline timing. Even when the source is a dynamic, evolving local outbreak rather than a single static node, the macroscopic simulation must correctly align its expected transmission time with the strict mathematical reality established at this microscopic baseline.


\section{Benchmark Case II: Chain of Cliques}
\label{sec:benchmark_chain_cliques}

\subsection{Construction of the Chain-of-Cliques Network}
\ReadStatusTwo{Construction of the Chain-of-Cliques Network}{Yulian}{1}{Michael}{0}

The second benchmark extends the single-export setting to a multi-community propagation problem. We consider a chain of \(M\) communities, where each community is a complete graph \(K_N\). Consecutive cliques are connected by sparse inter-community bridge links, while non-neighboring cliques are not directly connected. In the representative experiment shown below, the chain consists of \(M=5\) cliques with \(N=200\) nodes in each clique.

Formally, the node set is partitioned as
\[
V=\bigsqcup_{k=0}^{M-1}V_k,
\qquad
G[V_k]\simeq K_N,
\]
and inter-community edges are present only between neighboring communities,
\[
E_{k,k+1}\neq \varnothing,
\qquad
E_{ij}=\varnothing \quad \text{if } |i-j|>1.
\]
The process is initialized in the first community, and the main object of interest is the successive activation of communities along the chain.

This benchmark is a natural continuation of the clique-to-leaf case. The previous benchmark isolated one export event from a dense community. The present benchmark asks whether such export mechanisms remain accurate when they are composed repeatedly across several dense communities. It therefore tests not only local export behavior, but also the possible accumulation of timing error as the process moves farther from the initially active community.

\subsection{Expected Propagation Pattern}
\ReadStatusTwo{Expected Propagation Pattern}{Yulian}{1}{Michael}{0}

The chain-of-cliques benchmark is structurally favorable to the two-scale abstraction because each community is internally dense. Once a clique receives an imported infection, the local CTMC has many intra-community transmission opportunities, and the infection can spread rapidly through that clique. As a result, the global infected count is expected to have a step-like shape: relatively flat intervals correspond to waiting for the next successful inter-community export, while steep increases correspond to rapid within-clique spread after import.

The relevant comparison is therefore not only whether both simulators produce the same final infected count, but whether they reproduce the timing of the steps. Let \(T_k\) denote the first arrival time in community \(k\), with \(T_0\) determined by the initialization. Propagation along the chain can be decomposed as
\[
T_{k+1}=T_k+\Delta_k,
\]
where \(\Delta_k\) is the delay between the first activation of community \(k\) and the first import into community \(k+1\). This identity is only a temporal decomposition, not an independence assumption. Its importance is that small systematic differences in the delays \(\Delta_k\) can become more visible in downstream communities.

In the full-network reference, each delay \(\Delta_k\) is determined by the microscopic bridge-level state between two neighboring cliques. In the two-scale simulator, this bridge-level information is represented through the aggregate macro-rate
\[
\lambda_{k,k+1}(t)
=
\beta\, w_{k,k+1}\, C_k(t)\, S_{k+1}(t).
\]
Thus, this benchmark tests whether the macro summaries capture the effective timing of repeated inter-community transmissions in a setting where the local community structure is dense and relatively homogeneous.

\subsection{Numerical Comparison of Aggregate Infection Curves}
\ReadStatusTwo{Numerical Comparison of Aggregate Infection Curves}{Yulian}{1}{Michael}{0}

We first compare the full microscopic Gillespie simulator and the two-scale simulator using the aggregate infected count
\[
I(t)=\sum_{k=0}^{M-1} I_k(t).
\]
This observable summarizes the global progress of the process along the chain. Since each clique has \(N=200\) nodes, each major increase in \(I(t)\) corresponds to the activation and subsequent spread within one community.

\begin{figure}[t]
    \centering
    \includegraphics[width=\textwidth]{Micro_vs_MicroMacro.png}
    \caption{
    Aggregate infected trajectories for the chain-of-cliques benchmark with \(M=5\) cliques and \(N=200\) nodes per clique. 
    The figure compares the full microscopic Gillespie simulator and the two-scale simulator using representative trajectories, median curves, and interquartile bands. 
    The step-like shape reflects sequential activation of communities along the chain.
    }
    \label{fig:chain_cliques_aggregate_curves}
\end{figure}

Figure~\ref{fig:chain_cliques_aggregate_curves} shows that both simulators reproduce the qualitative stepwise propagation pattern. The aggregate infected count grows in several stages, consistent with sequential activation of the five cliques. The main difference is not the existence of these stages, but their timing and variability. This is the intended focus of the benchmark: in a dense chain of communities, the local outbreak shape is relatively regular, while the inter-community arrival times control the horizontal shifts between successive stages.

The interquartile bands are important because the comparison is distributional. Agreement of one representative trajectory would not be sufficient: the two-scale simulator should also reproduce the typical range of stochastic variation generated by the full-network reference. The aggregate plot therefore provides a first indication of whether the macro abstraction preserves the global propagation pattern and its uncertainty.

\subsection{Per-Community Activation Trajectories}
\ReadStatusTwo{Per-Community Activation Trajectories}{Yulian}{1}{Michael}{0}

To localize the aggregate behavior, we next compare representative per-community trajectories. For each community \(k\), the curve \(I_k(t)\) records the number of infected nodes in that clique. In this benchmark, the local growth after import is steep in both simulators, so the most informative quantity is the time at which each community first becomes activated.

\begin{figure}[t]
    \centering
    \includegraphics[width=\textwidth]{Micro_vs_MicroMacro_metric_panels.png}
    \caption{
    Representative per-community activation trajectories for the chain-of-cliques benchmark. 
    The left panel shows the full microscopic Gillespie simulation, and the right panel shows the two-scale simulation. 
    The curves show \(I_k(t)\) for each community, and the markers indicate selected activation-related times. 
    The comparison makes visible how differences in inter-community arrival times appear along the chain.
    }
    \label{fig:chain_cliques_per_community}
\end{figure}

Figure~\ref{fig:chain_cliques_per_community} shows that both simulators generate the same qualitative activation order imposed by the chain structure. Community \(0\) activates first, followed by the neighboring communities in sequence. After activation, each community undergoes a rapid local increase in \(I_k(t)\), which is consistent with the complete-graph structure inside each community.

At the same time, the activation times are not identical. The per-community view shows where timing differences enter: they appear primarily as shifts in the arrival times of later communities rather than as fundamentally different local growth shapes after import. This supports the interpretation that, in this benchmark, the dominant approximation effect is associated with repeated inter-community transmission rather than with the local micro-layer dynamics inside a clique.

\subsection{Interpretation of Error Accumulation}
\ReadStatusTwo{Interpretation of Error Accumulation}{Yulian}{1}{Michael}{0}

The chain-of-cliques benchmark reveals a mechanism that is not visible in a single clique-to-leaf export: approximation error can be composed over several inter-community transitions. In the full-network reference, each transition from community \(k\) to community \(k+1\) depends on the actual bridge-level configuration and on the current microscopic state near the boundary. In the two-scale simulator, this information is compressed into the macro-rate \(\lambda_{k,k+1}(t)\). A small mismatch in one export delay may shift the state from which the next export delay begins.

This effect is most relevant for downstream communities. If the first import into community \(k\) is advanced or delayed, then the local trajectory \(I_k(t)\) is also shifted. Since community \(k\) later becomes the source for community \(k+1\), the timing shift can propagate further along the chain. The relation
\[
T_{k+1}=T_k+\Delta_k
\]
therefore explains why later communities may show larger timing discrepancies even when each local clique behaves regularly after import.

In terms of the error mechanisms introduced earlier, this benchmark mainly probes rate abstraction under repeated composition. Import placement can also matter if the bridge links are concentrated on special boundary nodes, but its effect is less pronounced than in sparse or highly heterogeneous communities. Synchronization and numerical timing error should remain secondary provided that hazard accumulation and event-time interpolation are implemented consistently. Overall, the chain of cliques serves as an intermediate benchmark: it is more demanding than a single export motif, but still favorable to the two-scale approximation because the communities are dense and internally homogeneous.



\section{The Microscopic Ground Truth: The Clique and the Bridges}

\paragraph{Scaling Up: From a Single Node to a Community}
Studying the process of infection transmission  between two communities (one is infected, another is susceptible) we do not need the whole structure of target community but only it's nodes that are directly connected to the source community.
So, we now scale our perspective to model the interaction between an entire community and an individual. Instead of a single Infected source node, let us consider a source community modeled as a perfectly well-mixed, fully connected network—a \textit{clique}. This clique consists of a total of $n$ nodes. As time progresses, an internal epidemic spreads within this community, and the number of Infected nodes at any given time $t$ is denoted as $m(t)$. 

The target remains our isolated, Susceptible leaf node. However, this leaf is not connected to every single person in the source community. Instead, it is connected to a specific, predefined subset of exactly $k$ nodes within the clique. We will refer to these $k$ specific individuals as "bridge nodes," as they form the only physical pathways through which the infection can exit the community and reach the leaf.



\paragraph{The Phenomenon of "Local Blindness"}
To understand the precise timing of the transmission to the leaf, we must examine how the exact Gillespie algorithm processes this network topology. In a strict microscopic simulation, the infection pressure exerted on a node depends exclusively on the exact states of its immediate neighbors. 

This introduces a fundamental property of microscopic networks that we term \textit{Local Blindness}. From the perspective of the Susceptible leaf, the overall severity of the epidemic inside the clique is entirely invisible. The number of Infected nodes $m(t)$ could grow to encompass almost the entire community—say, $n-k$ nodes—but as long as the specific $k$ bridge nodes remain Susceptible, the leaf experiences absolutely zero infection risk. Its instantaneous hazard rate is mathematically strictly zero. The leaf only "feels" the epidemic, and its hazard integral only begins to accumulate, at the exact moment when the internal infection finally strikes one of the $k$ bridge nodes.

\paragraph{Deconstructing the Expected Transmission Time}
Because of this Local Blindness, the total time it takes for the leaf to become Infected in an exact microscopic simulation—let us call this $T_{micro}$—is fundamentally a two-step process. 

First, there is a delay phase. The infection must spread through the clique until it successfully "finds" and infects the first bridge node. The duration of this phase is highly dependent on the internal dynamics of the clique and the ratio of bridges $k$ to the total population $n$. Let us denote this time required to activate the first bridge as $T_{find}$. During this entire period, the risk to the leaf is zero.

Second, once the first bridge node transitions to the Infected state, an edge to the leaf becomes active. The process now mathematically mirrors the basic single-edge scenario we established in Section 2. We know that the expected time for the infection to cross a single active edge is $1/\beta$. Therefore, the expected total transmission time in the microscopic model can be conceptually understood as the expected time to find the bridge, plus the baseline expected time to cross it:
$$ \mathbb{E}[T_{micro}] \approx \mathbb{E}[T_{find}] + \frac{1}{\beta} $$
While the subsequent infection of additional bridge nodes would further increase the transmission rate and slightly accelerate the crossing phase, this additive two-step structure remains the defining characteristic of the exact simulation. The transmission is delayed, waiting for a specific topological event.

\paragraph{Establishing the Ground Truth for the Macro-Simulator}
This step-by-step, discrete process is characterized by periods of absolute safety followed by sudden exposure. Because the sequence of internal infections is random, the exact time $T_{micro}$ possesses significant variance. If a bridge node is infected early, the transmission occurs quickly; if the bridge nodes miraculously survive until the end of the internal outbreak, the transmission is vastly delayed.

This rigorous, path-dependent physical reality is our \textit{ground truth}. When we build a two-scale simulator and compress the entire source community into a single macroscopic node, we deliberately discard the ability to track these specific $k$ bridge nodes. Yet, if the two-scale simulator is to be scientifically valid, the macroscopic mathematical formula we use to trigger the inter-community jump must perfectly preserve and replicate this exact expected transmission time, $\mathbb{E}[T_{micro}]$. Understanding this ground truth is the prerequisite for identifying where naive macroscopic approximations fail.


\section{The Naive Macroscopic Approximation: A Mean-Field Approach}

\paragraph{The Need for Abstraction in Two-Scale Models}
As established in the previous section, the microscopic ground truth requires tracking the exact state of specific bridge nodes to determine when the infection jumps from a community to an external Susceptible leaf. However, the primary motivation for building a two-scale simulator is to escape the heavy computational burden of tracking individual nodes and edges. 

When we transition to the macro-layer of a two-scale simulator, we collapse the entire source community (the clique of size $n$) into a single macroscopic "metanode." In doing so, we deliberately discard the high-resolution topological map. The simulator no longer knows the identities of the $k$ bridge nodes, nor does it know precisely which individuals are currently in the \textit{Infected} state. The only information passed from the micro-layer to the macro-layer is a single, dynamically changing scalar value: $m(t)$, the total number of Infected nodes inside the community at time $t$. The challenge is how to calculate the transmission risk across the macroscopic edge using only this single aggregate number.



\paragraph{The Mean-Field Rationale: Probability as a Substitute for Topology}
To solve this, standard large-scale simulators often rely on a well-known mathematical shortcut called the \textit{mean-field approximation}. The logic behind this approach is highly intuitive. If we do not know exactly which nodes are Infected, we must treat all nodes in the community as mathematically identical.

We know that at any given time $t$, there are $m(t)$ Infected individuals out of a total population of $n$. If we assume "perfect mixing"—meaning the infection is distributed uniformly at random throughout the clique—what is the probability that any specific individual is Infected? It is simply the global fraction of Infected nodes, $m(t)/n$.

Because there are exactly $k$ bridge nodes connecting the community to the external Susceptible leaf, we apply this probability to the bridges. Instead of asking "how many bridge nodes are physically Infected right now?", the naive macroscopic approach asks "what is the \textit{expected} number of Infected bridge nodes?". By basic probability, the expected number of active bridges is the total number of bridges multiplied by the probability that any single one of them is Infected: $k \cdot \frac{m(t)}{n}$.

\paragraph{The Naive Macroscopic Hazard Function}
We can now translate this expected value into a mathematical function for the two-scale simulator. In the fundamental single-edge baseline (Section 2), the transmission rate across one active edge is $\beta$. If we expect to have $k \cdot \frac{m(t)}{n}$ active edges, we simply multiply this expected number by the baseline rate.

This gives us the naive macroscopic hazard function, $\lambda_{macro}(t)$, which calculates the continuous infection pressure exerted by the entire community onto the Susceptible leaf:
$$ \lambda_{macro}(t) = k \cdot \frac{m(t)}{n} \cdot \beta $$

In this formulation, the instantaneous risk of the infection crossing the macro-edge is directly proportional to the overall prevalence of the disease inside the source community. As the local outbreak grows (as $m(t)$ increases), the pressure on the external leaf smoothly and proportionally rises.

\paragraph{The Computational Advantage}
It is easy to see why this formula is universally appealing for large-scale epidemic architectures. It is extremely fast and computationally lightweight. 

To evaluate whether the infection jumps to the next community, the macro-simulator only needs to read the scalar value $m(t)$ from the micro-layer, plug it into the simple arithmetic formula above, and integrate the result. It completely bypasses the need to maintain an $n \times n$ adjacency matrix or iterate through the complex micro-topology of the bridges. By replacing strict deterministic pathways with a smooth, probability-driven average, the simulator achieves massive scalability. However, as we will demonstrate in the next section, substituting exact topological states with this smooth macroscopic average introduces a hidden and severe temporal flaw.



\section{The Timing Discrepancy: Premature Leakage and Systemic Bias}

\paragraph{The Illusion of the Mean-Field Average}
In the previous section, we established that the naive macroscopic approximation elegantly solves the computational problem by assuming the infection pressure is proportional to the global fraction of Infected nodes, $m(t)/n$. On the surface, this mean-field approach appears mathematically sound: over a large number of independent trials, the average number of Infected bridge nodes indeed approaches $k \cdot m(t)/n$. However, substituting an exact topological requirement with a probabilistic average introduces a fatal flaw in the temporal dynamics of the simulation.

To understand this flaw, we must contrast how the "risk clock" operates in the two models. In continuous-time stochastic simulations, transmission occurs when the accumulated risk (the hazard integral) reaches a randomly drawn threshold. The time it takes to accumulate this risk strictly determines the speed of the epidemic.

\paragraph{The Concept of "Premature Leakage"}
The critical discrepancy lies in the exact moment the hazard integral begins to accumulate values greater than zero. 

Recall the microscopic ground truth (Section 3) and the phenomenon of "Local Blindness." In the exact physical reality, the Susceptible leaf is completely shielded from the internal epidemic as long as the infection remains among non-bridge nodes. The risk clock strictly reads zero, and time essentially stands still for the leaf until the specific moment a bridge node becomes Infected. 

Now consider the naive macroscopic hazard function: $\lambda_{macro}(t) = k \cdot \frac{m(t)}{n} \cdot \beta$. 
The moment the very first node inside the source community becomes Infected (i.e., $m(t) = 1$), the fraction $m(t)/n$ becomes greater than zero. Consequently, the macroscopic hazard rate immediately becomes positive. The risk clock starts ticking, and the hazard integral begins to accumulate risk from the very beginning of the local outbreak. 

We term this phenomenon \textit{Premature Leakage}. The naive model allows transmission risk to "leak" out of the community and exert pressure on the external Susceptible leaf long before the physical infection has actually navigated the internal topology to reach a valid bridge node. 



\paragraph{The Mathematical Consequence: Underestimating Transmission Time}
Because the naive macroscopic model starts integrating risk prematurely, the area under its hazard curve grows much faster in the early stages of the local outbreak compared to the exact microscopic model. 

Mathematically, if a function begins accumulating area earlier, it will inevitably reach the required stochastic threshold, $\theta = -\ln(U)$, sooner. While the exact model waits patiently during the delay phase ($T_{find}$) before steeply accumulating risk, the naive model continuously chips away at the threshold from day one. 

As a direct mathematical consequence, the naive macroscopic formula systematically underestimates the time required for the infection to cross between communities. If we run many simulations, the expected time generated by the macroscopic approximation will always be strictly less than the true expected time generated by the exact microscopic model:
$$ \mathbb{E}[T_{macro}] < \mathbb{E}[T_{micro}] $$
This is not merely a random fluctuation or stochastic noise; it is a persistent, systemic bias mathematically hardcoded into the mean-field approximation.

\paragraph{Severe Implications for the Two-Scale Simulator}
The presence of this systemic bias has devastating consequences for the reliability of a large-scale, two-scale simulator. 

Epidemic modeling on a national or global scale relies heavily on the accurate cascading of infections from one community to the next. If the hazard function used for the macroscopic edges suffers from Premature Leakage, the infection will jump from the first community to the second artificially fast. When the second community experiences its internal outbreak, it will prematurely leak risk to a third community, and so on.

Because $\mathbb{E}[T_{macro}] < \mathbb{E}[T_{micro}]$, these temporal errors are not isolated—they compound exponentially as the disease traverses the macroscopic network. A two-scale simulator running this naive approximation will behave like a fast-forwarded movie. It will predict that the epidemic spreads across regions much faster than physically possible, artificially accelerating the global timeline. Consequently, public health interventions based on such a model would be mistimed, as the simulation would project the peak of the pandemic weeks or even months earlier than the true topological dynamics dictate. Therefore, to make the two-scale architecture viable, this premature leakage must be mathematically corrected.



\subsubsection{Experiment on macro-clique of complete graphs}
number of communities = 5
number of nodes inside each community = 100
number of links between comminities = 5

Topology of communities = complete graph
Topology of MacroGraph = complete graph


\begin{figure}[hbt!]
    \centering

    % --- First Row (2 Images) ---
    \begin{subfigure}[b]{0.48\textwidth}
        \centering
        % Filenames with spaces are best wrapped in quotes or handled by modern compilers
        \includegraphics[width=\linewidth]{Images/updated_network_k5_n100_inter5_macrocomplete_microcomplete_p0p2/Micro vs MicroMacro.png}
        \caption{Micro vs MicroMacro}
        \label{fig:micro_vs_micromacro_updated}
    \end{subfigure}
    \hfill
    \begin{subfigure}[b]{0.48\textwidth}
        \centering
        \includegraphics[width=\linewidth]{Images/updated_network_k5_n100_inter5_macrocomplete_microcomplete_p0p2/Micro_community.png}
        \caption{Micro Community}
        \label{fig:micro_community_updated}
    \end{subfigure}

    \bigskip % Vertical space between rows

    % --- Second Row (1 Image Centered) ---
    \begin{subfigure}[b]{0.48\textwidth}
        \centering
        \includegraphics[width=\linewidth]{Images/updated_network_k5_n100_inter5_macrocomplete_microcomplete_p0p2/MicroMacro_community.png}
        \caption{MicroMacro Community}
        \label{fig:micromacro_community_updated}
    \end{subfigure}

    % Main Caption
    \caption{Gallery for updated\_network\_k5\_n100\_inter5\_macrocomplete\_microcomplete\_p0p2}
    \label{fig:updated_network_gallery}
\end{figure}

\subsubsection{Chain of communities}
number of communities = 5
number of nodes inside each community = 100
number of links between communities = 5

Topology of communities = complete graph
Topology of MacroGraph = complete graph

\begin{figure}[hbt!]
    \centering

    % --- First Row ---
    \begin{subfigure}[b]{0.48\textwidth}
        \centering
        \includegraphics[width=\linewidth]{Images/k5_n100_inter5_macrochain_microcomplete_p0p2/network_k5_n100_inter5_macrochain_microcomplete_p0p2.png}
        \caption{Network Visualization}
        \label{fig:macrochain_network}
    \end{subfigure}
    \hfill
    \begin{subfigure}[b]{0.48\textwidth}
        \centering
        \includegraphics[width=\linewidth]{Images/k5_n100_inter5_macrochain_microcomplete_p0p2/MicroMacro_vs_Micro.png}
        \caption{MicroMacro vs Micro}
        \label{fig:macrochain_micromacro_vs_micro}
    \end{subfigure}

    \bigskip % Vertical space between rows

    % --- Second Row ---
    \begin{subfigure}[b]{0.48\textwidth}
        \centering
        \includegraphics[width=\linewidth]{Images/k5_n100_inter5_macrochain_microcomplete_p0p2/MicroMacro_community.png}
        \caption{MicroMacro Community}
        \label{fig:macrochain_micromacro_community}
    \end{subfigure}
    \hfill
    \begin{subfigure}[b]{0.48\textwidth}
        \centering
        \includegraphics[width=\linewidth]{Images/k5_n100_inter5_macrochain_microcomplete_p0p2/Micro_community.png}
        \caption{Micro Community}
        \label{fig:macrochain_micro_community}
    \end{subfigure}

    % Main Caption
    \caption{Chain of cliques}
    \label{fig:macrochain_gallery}
\end{figure}

\begin{figure}[hbt!]
    \centering

    % --- First Row ---
    \begin{subfigure}[b]{0.48\textwidth}
        \centering
        % Matches the network visualization file
        \includegraphics[width=\linewidth]{Images/k_5_n_100_inter_5_macro_chain_micro_random_p0p2/network_k5_n100_inter5_macrochain_microrandom_p0p2.png}
        \caption{Network Visualization}
        \label{fig:microrandom_network}
    \end{subfigure}
    \hfill
    \begin{subfigure}[b]{0.48\textwidth}
        \centering
        % Matches the comparison plot
        \includegraphics[width=\linewidth]{Images/k_5_n_100_inter_5_macro_chain_micro_random_p0p2/Micro_vs_MicroMacro_v2.png}
        \caption{Micro vs MicroMacro (v2)}
        \label{fig:microrandom_micromacro_vs_micro}
    \end{subfigure}

    \bigskip % Vertical space between rows

    % --- Second Row ---
    \begin{subfigure}[b]{0.48\textwidth}
        \centering
        % Matches the MicroMacro community plot
        \includegraphics[width=\linewidth]{Images/k_5_n_100_inter_5_macro_chain_micro_random_p0p2/MicroMacro_v2_community.png}
        \caption{MicroMacro Community}
        \label{fig:microrandom_micromacro_community}
    \end{subfigure}
    \hfill
    \begin{subfigure}[b]{0.48\textwidth}
        \centering
        % Matches the Community comparison plot (closest match to "Micro Community")
        \includegraphics[width=\linewidth]{Images/k_5_n_100_inter_5_macro_chain_micro_random_p0p2/Micro_vs_MicroMacro_v2_community.png}
        \caption{Micro vs MicroMacro Community}
        \label{fig:microrandom_micro_community_compare}
    \end{subfigure}

    % Main Caption
    \caption{Chain of cliques (Micro Random, $p=0.2$)}
    \label{fig:microrandom_gallery}
\end{figure}



\section{Conclusions}

The results suggest that the two-scale approximation is most reliable when boundary nodes are representative of their communities, when inter-community interfaces are sufficiently broad, and when export events are rare relative to within-community mixing. 
In contrast, narrow bridge bottlenecks, strong correlations between infection status and boundary position, and high export rates can produce systematic premature leakage. 
Thus, the macro layer should not be interpreted as an exact replacement for the full inter-community edge structure, but as an effective approximation whose validity depends on the relation between community mixing and boundary exposure.



\bibliographystyle{apalike}
\bibliography{references}

\end{document}
